\documentclass[11pt,a4paper]{article}
\usepackage[T1]{fontenc}
\usepackage{lmodern}
\usepackage[margin=25mm,headheight=15pt]{geometry}
\usepackage{amsmath,amssymb,amsthm,mathtools}
\usepackage{booktabs,longtable,array,enumitem}
\newcolumntype{P}[1]{>{\raggedright\arraybackslash}p{#1}}
\usepackage{microtype,xcolor,xurl}
\usepackage[hidelinks]{hyperref}
\usepackage{fancyhdr}
\usepackage{listings}
\newcommand{\R}{\mathbb{R}}
\newcommand{\E}{\mathbb{E}}
\newcommand{\Ind}{\mathbf{1}}
\DeclareMathOperator{\clip}{clip}
\DeclareMathOperator{\softplus}{softplus}
\DeclareMathOperator{\Var}{Var}
\DeclareMathOperator{\Cov}{Cov}
\DeclareMathOperator{\median}{median}
\DeclareMathOperator{\diag}{diag}
\DeclareMathOperator{\argmin}{argmin}
\newtheorem{proposition}{Proposition}[section]
\newtheorem{remark}{Remark}[section]
\setlist{itemsep=3pt,topsep=5pt}
\setlength{\emergencystretch}{2em}
\setlength{\LTpre}{7pt}
\setlength{\LTpost}{7pt}
\pagestyle{fancy}
\fancyhf{}
\fancyhead[L]{\small WildfireGuardian: mathematical methods}
\fancyhead[R]{\small 8 October 2026}
\fancyfoot[C]{\thepage}
\lstset{basicstyle=\ttfamily\footnotesize,breaklines=true,columns=fullflexible,
 keepspaces=true,showstringspaces=false,frame=single,rulecolor=\color{black!20},
 aboveskip=9pt,belowskip=9pt}
\hypersetup{pdftitle={WildfireGuardian: mathematical methods for burned-area prediction, active flames, and expected radiant heat},pdfauthor={WildfireGuardian methods documentation}}
\title{\textbf{WildfireGuardian: Mathematical Methods}\\[5pt]
\large A step-by-step account of retrospective arrival reconstruction,\\
hourly transition prediction, flame residence, and expected radiant heat}
\author{Methods document grounded in the current project implementation}
\date{8 October 2026 (Asia/Seoul)}
\begin{document}
\maketitle
\begin{abstract}
This document specifies the maintained WildfireGuardian workflow as it exists
on 8 October 2026. It follows the calculation from published satellite-based
time-of-arrival reconstructions to aligned environmental tensors, two trained
environment-conditioned spatial models, assumed active-flame probabilities,
and an expected incident radiant-heat field. Each mathematical object is tied
to its units, spatial support, time convention, validity conditions, and
implementation. The nonnegative exponential-link model and signed direct-score
model are treated separately. Derivations distinguish reconstructed labels,
learned next-hour transition probabilities, deterministic residence scenarios,
and the additional physical and timing assumptions required for radiation.
Saved comparisons and demonstrations are reported as recorded evidence, not
as fresh model fitting or independent physical validation. In particular,
375 m raster spacing and hourly output do not imply corresponding observation
accuracy; a cumulative burned state is not an active-flame map; and a mean
radiant flux does not determine a heat-exceedance probability.
\end{abstract}
\clearpage
\tableofcontents
\clearpage

\section{Scope, objects, and the calculation being performed}
\label{sec:scope}
\subsection{Exactly which workflow is described}
The scope is the current root \path{README.md}, the maintained Python packages
\path{wildfire_burned_area_hourly/} and \path{expected_heat_flux/}, their saved
experiments, and the publisher reconstruction code retained under
\path{wildfire_burned_area_hourly/sources/zenodo_21788344/extracted/code/}.
The published reconstruction is an upstream method. The local bundle consumes
its supplied rasters and constructs labels from them; ordinary model inference
does not re-run kriging. Other historical projects, routing algorithms, and
previous neural-network packages are outside this document's scope.

There are six successive operations:
\begin{enumerate}
\item Estimate an arrival time at each covered location from event-wide
active-fire detections and a known final perimeter, using the published
variance-weighted ordinary-kriging procedure.
\item Compare arrival time with each snapshot's event-relative clock to obtain
cumulative burned-state and hourly transition labels.
\item Align environmental sources to the original rectangular grid, retaining
each source's coverage, units, native support, and categorical vocabulary.
\item Fit or load an environment-conditioned spatial model for a currently
unburned target cell's first new burning during the next hour.
\item Supply a flame or burning duration and an arrival-time assumption to
convert that transition probability into activity at an instant or mean activity
over an interval.
\item Supply dry fuel, combustion, radiation, and geometry assumptions to
integrate the activity-weighted emitting sources into a radiant-heat field.
\end{enumerate}

Every later operation depends on the earlier contracts, but it does not erase
their uncertainty or turn retrospective information into live observations.
The learned probability model receives only current-state and current
environmental tensors at inference. Nevertheless, the current burned-state
tensor itself was reconstructed using final extent and event-wide information.
Restricting tensor timestamps to the cutoff does not make that input
operationally available at the cutoff.

\subsection{Global notation and units}
\label{sec:notation}
The following notation is used throughout. Sections may introduce additional
locally defined indices; the domain and units are stated where introduced.
\begin{longtable}{@{}P{0.17\linewidth}P{0.76\linewidth}@{}}
\toprule Symbol & Definition and units\\\midrule
\endfirsthead
\toprule Symbol & Definition and units\\\midrule
\endhead
$e$ & A fire event with its own raster, coordinate reference system, and time origin.\\
$\mathcal G_e$ & The finite rectangular grid of event $e$; a cell index is $i=(r,c)$ with zero-based row $r$ and column $c$.\\
$S_i$ & The horizontal ground footprint of cell $i$, in a projected metric plane; its area is $|S_i|$ in $\mathrm{m^2}$.\\
$x_i$ & The projected horizontal center of $S_i$, in metres.\\
$t_{0,e}$ & Event-specific UTC origin of the arrival raster.\\
$t$, $h_e(t)$ & Absolute UTC snapshot time and elapsed hours $h_e(t)=(t-t_{0,e})/(3600\ \mathrm{s})$.\\
$\widehat a_i$ & Reconstructed arrival time, in hours since $t_{0,e}$; defined only where its raster mask is valid.\\
$M_i^A$ & Boolean arrival-coverage indicator; one denotes a known reconstruction, zero denotes unavailable arrival.\\
$B_i(t)$ & Reconstructed cumulative burned state, $\Ind\{\widehat a_i\le h_e(t)\}$ where $M_i^A=1$; dimensionless.\\
$Y_i(t)$ & Binary label for new burning during $(t,t+1\ \mathrm{h}]$, on a known currently unburned target.\\
$E_i(t)$ & Ordered environmental feature vector; raw continuous entries retain their recorded physical units and one-hot entries are dimensionless.\\
$s_i(t)$ & Unclipped signed direct-model score, dimensionless; it may lie outside $[0,1]$.\\
$\lambda_i(t)$ & Nonnegative exponential-model rate in $\mathrm{h^{-1}}$, under a frozen current input state.\\
$q_i(t)$ & Usable next-hour new-burn probability, dimensionless in $[0,1]$, or unavailable with a false prediction mask.\\
$D$ & Duration of the heat forecast interval in seconds; the composed current model requires $D=3600\ \mathrm{s}$.\\
$T_i$ & Supplied positive conditional burning duration for the heat calculation, in seconds.\\
$\tau_i^{\mathrm{fl}}$ & Surface-flame residence time in seconds; it equals $T_i$ only when explicitly made consistent by the caller.\\
$p_i(u)$, $\overline p_i$ & Active-burning probability at elapsed seconds $u$ and its time average over $[0,D]$; dimensionless.\\
$w_i$, $f_i$, $H_{c,i}$ & Dry fuel loading in $\mathrm{kg\,m^{-2}}$, consumed fraction in $[0,1]$, and effective heat of combustion in $\mathrm{J\,kg^{-1}}$.\\
$Q_i$ & Heat-release-rate density conditional on modeled activity, in $\mathrm{W\,m^{-2}}$.\\
$\chi_i$ & Radiative fraction of conditional heat release, dimensionless in $[0,1]$.\\
$\mathcal T$ & Atmospheric transmissivity, dimensionless in $[0,1]$; distinct from duration and from flame residence.\\
$\Omega_i(x)$ & Solid angle of elevated source footprint $S_i$ seen from receiver point $x$, in steradians.\\
$F(x,u)$ & Incident radiant flux at a horizontal point receiver, in $\mathrm{W\,m^{-2}}$; mean flux and endpoint flux are different outputs.\\
$\mathcal E(x)$ & Time-integrated expected incident radiant energy, in $\mathrm{J\,m^{-2}}$.\\
\bottomrule
\end{longtable}

The code field named \path{atmospheric_transmissivity} is denoted by
$\mathcal T$ here even when a later equation or quoted method uses another
letter. The field \path{burning_duration_s} is $T_i$; the residence object's
\path{duration_s} is $\tau_i^{\mathrm{fl}}$. A notation reused within a section
is explicitly identified. Tensor time length is denoted $L$, and batch size
$N_b$, to avoid confusing these array dimensions with physical duration or
burned state.

\subsection{Six statements that are not interchangeable}
\begin{enumerate}
\item A sensor detected active fire at a location and acquisition time.
\item A retrospective arrival reconstruction is valid at a raster cell.
\item A raster cell is classified as cumulatively burned at a snapshot time.
\item A model assigns probability $q_i$ to its first new burning in the next hour.
\item A timing-and-residence scenario assigns probability $p_i(u)$ to a
representative location being actively flaming at an instant.
\item A radiation scenario gives an expected incident flux at a receiver.
\end{enumerate}
The numerical outputs support only the corresponding statement. For example,
$B_i=1$ does not imply present activity, and $q_i=0.8$ does not imply an 80\%
chance of flames at the one-hour endpoint. A validity mask is required even
when a stored numeric value is zero.

\section{Step 1: published retrospective arrival-time reconstruction}
\label{sec:reconstruction}
\subsection{Source product and observation preprocessing}
The retained deposit supplies 30 Korean rasters (15 fires times arrival and
kriging standard deviation) and 58 U.S. rasters (29 fires times the same two
quantities). Korean output grids use EPSG:5179 and U.S. output grids use
EPSG:5070. Both have 375 m output spacing. The source is the dataset and code
identified by DOI \href{https://doi.org/10.5281/zenodo.21788344}{10.5281/zenodo.21788344}
\cite{toadeposit,codetoa}. This document explains the archived implementation;
it does not claim that the raw detections or aerial reference perimeters were
downloaded again or that the published study was independently reproduced.

Let an input detection record contain longitude, latitude, UTC acquisition time,
and sensor. The archived \path{preprocess.select} removes fully identical rows,
retains times within a closed collection interval $[t_{\min},t_{\max}]$, and
uses strict spatial inclusion in the final perimeter. Geostationary detections
are selected inside a perimeter expanded by 2000 m; VIIRS uses no expansion.
Boundary points are excluded by the implemented strict containment predicate.

U.S. analysis coordinates are EPSG:5070 coordinates divided by 1000, in km.
Korean analysis coordinates are the source's approximation
\begin{equation}
x=111\,\ell\cos(\varphi_0),\qquad y=111\,\varphi,
\end{equation}
where longitude $\ell$, latitude $\varphi$, and event-center latitude
$\varphi_0$ are in degrees and the outputs are in km. The cosine argument is
converted to radians in code. These Korean \emph{analysis} coordinates differ
from the projected EPSG:5179 coordinates of the delivered raster.

VIIRS detections are grouped using
\begin{equation}
(g_x,g_y)=\left(\operatorname{round}(x/0.375),
                     \operatorname{round}(y/0.375)\right).
\end{equation}
The implementation uses NumPy rounding, including its ties-to-even convention.
Each occupied analysis cell contributes the median detection time; its
representative location is the earliest detection's location, breaking time
ties by smaller latitude and then longitude. Thus the temporal statistic is
the \emph{median}, whereas the representative spatial point comes from the
\emph{earliest} detection. Neither operation establishes exact physical
ignition time.

\subsection{Event-relative clocks}
Each detection time is converted to hours after its event origin. In the U.S.,
the origin is midnight UTC on the day containing the start of the collection
window, whose start is the first NIROPS survey minus 24 hours. For Korea,
the origin is the Korea Forest Service situation-map ignition time.
The authoritative local origin comes from GeoTIFF tags and event metadata.
Never compare arrival numbers from two events before adding their respective
origins. Valid U.S. arrival predictions may be negative: validity is read from
NoData and coverage metadata, not inferred from the sign.

\subsection{Empirical semivariogram and fitted covariance}
Let $z_j$ be the detection-cell summary time in hours at analysis point
$\xi_j\in\R^2$ in km, for $j=1,\ldots,n$. For each nonempty lag bin
$\mathcal P_k$ of unordered point pairs, define
\begin{equation}
\widehat h_k=\frac1{|\mathcal P_k|}\sum_{(j,l)\in\mathcal P_k}
\|\xi_j-\xi_l\|_2,\qquad
\widehat\gamma_k=\frac1{2|\mathcal P_k|}
\sum_{(j,l)\in\mathcal P_k}(z_j-z_l)^2.
\end{equation}
The code uses 20 equal-width bins between the smallest and largest pair
distance, with a final edge extended by 0.001 km. Empty bins are omitted.
For more than 15,000 observations it accumulates the same pair statistics in
chunks instead of materializing all pairs at once.

The fitted model is
\begin{equation}
\gamma(h)=p\left(1-e^{-3h/r_0}\right)+c_0,
\qquad p\ge0,\quad r_0>0,\quad c_0\ge0,
\label{eq:variogram}
\end{equation}
where $p$ is the partial sill in $\mathrm{h^2}$, $c_0$ is the nugget in
$\mathrm{h^2}$, and $r_0$ is the effective range in km. At $h=r_0$ the
structured term is $(1-e^{-3})p\simeq0.9502p$. Bounded SciPy least squares
uses the \path{soft_l1} loss on residuals
$\gamma(\widehat h_k)-\widehat\gamma_k$ with its default unit scale, tight
termination tolerances $10^{-15}$, and at most 2,000,000 function evaluations.
The initial values are
\begin{equation}
\left(\max_k\widehat\gamma_k-\min_k\widehat\gamma_k,
       \tfrac14\max_k\widehat h_k,
       \min_k\widehat\gamma_k\right),
\end{equation}
and upper parameter bounds are
$(10\max_k\widehat\gamma_k,\max_k\widehat h_k,
\max_k\widehat\gamma_k)$. The returned nugget is floored to
$\max(c_0,0.001p)$. A failed convergence status raises an error.
The implemented covariance is
\begin{equation}
C(h)=\begin{cases}p+c_0,&h=0,\\p e^{-3h/r_0},&h>0.\end{cases}
\end{equation}
The nugget at zero is an explicit implementation convention, not a claim that
the semivariogram expression at zero directly equals the variance of a
difference between a variable and itself.

\subsection{Local ordinary kriging with observation-dependent error variance}
At a target point $\xi$, choose the nearest $m=\min(50,n_{\mathrm{all}})$
observations from the combined VIIRS and geostationary table. At least three
neighbors are required for a finite prediction. The production procedure also
requires at least five VIIRS summary cells to fit an event product.
The variogram is fitted on VIIRS summary cells, before adding geostationary
points to the kriging system.

Let $v_j$ denote assigned error variance, in $\mathrm{h^2}$:
\begin{equation}
v_j=\begin{cases}
0,&\text{VIIRS summary cell},\\
27.4\max(c_0,1\ \mathrm{h^2}),&\text{geostationary point}.
\end{cases}
\end{equation}
The source's factor 27.4 is a rounded resolution-based inflation factor; its
comment motivates it by $(2/0.375)^2-1$. It is a prescribed error model,
not an empirically calibrated uncertainty for every geostationary observation.
Write $\epsilon=10^{-4}p$ for the numerical ridge, and set
\begin{equation}
K_{jl}=C(\|\xi_j-\xi_l\|_2)+\Ind\{j=l\}(v_j+\epsilon),
\qquad c_j=C(\|\xi_j-\xi\|_2).
\end{equation}
The weights $\alpha\in\R^m$ and Lagrange multiplier $\mu\in\R$ solve
\begin{equation}
\begin{pmatrix}K&\boldsymbol 1\\\boldsymbol 1^\top&0\end{pmatrix}
\begin{pmatrix}\alpha\\\mu\end{pmatrix}
=\begin{pmatrix}c\\1\end{pmatrix}.
\label{eq:krigingsystem}
\end{equation}
The last row enforces $\boldsymbol1^\top\alpha=1$, the ordinary-kriging
unbiasedness constraint under an unknown constant mean.
To derive the system, minimize the mean-squared error
\begin{equation}
V(\alpha)=C(0)-2\alpha^\top c+\alpha^\top K\alpha
\end{equation}
subject to that constraint. Differentiating the Lagrangian
$V(\alpha)+2\mu(\boldsymbol1^\top\alpha-1)$ yields
$K\alpha+\mu\boldsymbol1=c$, exactly the first block row.
Therefore
\begin{equation}
\widehat a(\xi)=\alpha^\top z,\qquad
\widehat V(\xi)=\max\{p+c_0-\alpha^\top c-\mu,0\}.
\end{equation}
The variance simplification follows by multiplying the first block row by
$\alpha^\top$: $\alpha^\top K\alpha=\alpha^\top c-\mu$.
The nonnegative floor is a numerical safeguard. The delivered uncertainty
raster is $\sqrt{\widehat V}$ in hours. It is a model-based kriging standard
deviation; the local downstream methods do not treat it as a calibrated
Gaussian posterior, sample arrival-time realizations, or convert it into
flame-confidence intervals.

\subsection{Reference evaluation and internal cross-validation}
The upstream code compares five interpolation configurations: VIIRS-only IDW,
VIIRS-plus-geostationary IDW, VIIRS-only ordinary kriging, pooled unweighted
ordinary kriging, and the variance-weighted pooled method above. IDW uses
20 nearest observations and weights proportional to
$\max(\|\xi_j-\xi\|_2,10^{-10}\ \mathrm{km})^{-2}$.
These are reconstruction comparators, distinct from the later learned
exponential and signed direct probability models.

U.S. external reference arrival is the time of the first aerial infrared survey
whose cumulative union encloses a cell center. It is stepwise at survey times,
not continuously observed physical arrival. At survey time $t_s$, let $H_s$,
$M_s$, and $F_s$ be counts of intersection, reference-only, and prediction-only
burned cells, respectively, within the final reference support. The implemented
S\o rensen coefficient is
\begin{equation}
\mathrm{SC}_s=\frac{2H_s}{2H_s+M_s+F_s}.
\end{equation}
The default growth window retains surveys with reference count
$\max(1,0.01N)\le N_s<0.90N$, where $N$ is the final reference count. An event
score is the median over eligible surveys, rounded to four decimals.

Spatial-block internal validation withholds VIIRS cells while retaining
geostationary observations. Block width is
\begin{equation}
b=\min\{30,\max(5,r_0/2)\}\quad\mathrm{km}.
\end{equation}
The U.S. procedure withholds each occupied block, requires at least five
training cells, refits the variogram for each fold, and uses a fold only when
all five configurations return finite predictions. The Korean procedure uses
at most ten grouped folds, at least 20 training and three validation cells
per fold, and adoption conditions including whole-data range at least 0.5 km,
three occupied blocks, ten pooled validation pairs, and nonzero observation
and prediction variance. Withheld evaluation is an unavailable assessment,
not proof of inaccurate reconstruction.

For residuals $d_j=\widehat z_j-z_j$, timing scores are
\begin{equation}
\mathrm{ME}=\frac1n\sum_jd_j,\qquad
\mathrm{MAE}=\frac1n\sum_j|d_j|,
\end{equation}
in hours, with Pearson correlation also reported. Across events, paired exact
two-sided sign tests omit zero differences, and Holm adjustment controls the
family of ten configuration pairs. If ordered raw values are $p_{(1)}\le
\cdots\le p_{(10)}$, adjusted values are
$\widetilde p_{(k)}=\max_{j\le k}\min\{1,(11-j)p_{(j)}\}$.
Percentile intervals for median paired differences and Spearman correlations
resample entire event pairs 10,000 times with saved seed 20260906.
These source-study statistical procedures were read from retained code;
they were not rerun to author this document.

\section{Step 2: converting arrival time to hourly labels}
\label{sec:labels}
\subsection{Cumulative state and transition intervals}
For $M_i^A=1$, the reconstructed cumulative state is
\begin{equation}
B_i(t)=\Ind\{\widehat a_i\le h_e(t)\}.
\end{equation}
It is nondecreasing in $t$. At the precise reconstructed arrival time it is
already one. Where $M_i^A=0$, the physical state remains unknown even though a
storage tensor may contain zero.
The loader's field \path{newly_burned_area} at snapshot $t$ means the
\emph{preceding} hour:
\begin{equation}
N_i^{\mathrm{past}}(t)=
\Ind\{h_e(t)-1<\widehat a_i\le h_e(t)\}.
\label{eq:pastlabel}
\end{equation}
The model's training label instead means the \emph{following} hour:
\begin{equation}
Y_i(t)=\Ind\{h_e(t)<\widehat a_i\le h_e(t)+1\}
       =B_i(t+1\ \mathrm h)-B_i(t).
\label{eq:futurelabel}
\end{equation}
This equality assumes the same valid reconstructed arrival is used at both
times. Training candidates have known $B_i(t)=0$, so the future burned value
also equals $Y_i(t)$ on those candidates. The training preparation loads current environmental and burned-state
contexts in 24-hour chunks. For a chosen cutoff, it computes the next-hour
target state directly from the arrival raster using the next elapsed-hour
boundary, then checks equality with the independently enumerated candidate
label. It does not use the loader's preceding-hour field as the forecast
target. The precise eligibility and sampling procedure is given
in the probability-method section.

\subsection{Mask encoding, area, and temporal precision}
Exported categorical hourly masks use code 1 for estimated burned, code 0 for
estimated not yet burned \emph{inside} valid reconstruction support, and code
255 for unavailable. Code 0 outside a tensor's validity mask cannot be
interpreted as measured unburned ground.
At square 375 m spacing,
\begin{equation}
|S_i|=375^2=140625\ \mathrm{m^2}=14.0625\ \mathrm{ha}.
\end{equation}
The cumulative raster-support area is
$\mathcal A_e(t)=14.0625\sum_{i:M_i^A=1}B_i(t)$ hectares. It counts entire projected
supports selected by a center-in-perimeter rule. It need not equal the exact
vector-perimeter area or administrative forest-damage area, and no such area
equality or rescaling is imposed.

One-hour output spacing is a thresholding schedule for a continuous
reconstruction, not a one-hour accuracy statement. For example, the retained
Uljin regional metadata reports block-validation MAE about 10.44 hours and
kriging standard deviation about 4.40--52.96 hours. The current active-flame
example has much larger arrival uncertainty than its seconds-long residence
scenario. Thresholding uncertainty at a finer clock does not remove it.

\section{Step 3: aligning and loading environmental inputs}
\label{sec:environment}
\subsection{Spatial alignment and the original rectangular grid}
The environmental package retains each event's original arrival-raster shape,
CRS, and affine transform. The environment occupies the \emph{rectangular}
grid and is not cropped to future final burned extent. Arrival labels retain
their own reconstruction mask. Write the affine coefficients as
$(a,b,c,d,e,f)$, using the conventional raster corner transform. For zero-based
row $r$ and column $c_{\mathrm{col}}$, the cell center is
\begin{equation}
\begin{pmatrix}x_{r,c_{\mathrm{col}}}\\y_{r,c_{\mathrm{col}}}\end{pmatrix}
=\begin{pmatrix}a&b\\d&e\end{pmatrix}
\begin{pmatrix}c_{\mathrm{col}}+1/2\\r+1/2\end{pmatrix}
+\begin{pmatrix}c\\f\end{pmatrix}.
\end{equation}
The corner coordinate $c$ in this equation is an affine coefficient, not the
column index. The cell area is $|ae-bd|$ in square metres. On an extracted
patch beginning at row $r_0$, column $c_0$, its affine transform is the event
transform composed with pixel translation $(c_0,r_0)$. Coordinates and masks
are supplied along with values; padded spatial coordinates are NaN.

\subsection{Source support and resampling}
\begin{longtable}{@{}P{0.23\linewidth}P{0.31\linewidth}P{0.38\linewidth}@{}}
\toprule Input family & Source/support in this bundle & Alignment and interpretation\\\midrule
\endfirsthead
\toprule Input family & Source/support in this bundle & Alignment and interpretation\\\midrule
\endhead
Weather & Archived IFS model values, roughly 9 km and hourly & Nearest returned native model-node center in the event metric plane; all its assigned pixels share the same node value.\\
Mixing depth & ERA5, about $0.25^\circ$ and hourly & Nearest native-node mapping; boundary-layer height is a mixing-depth proxy.\\
Terrain & Copernicus GLO-30 DSM, nominal 30 m & Bilinear reprojection to 31.25 m intermediate grid, terrain derivatives there, then finite-value aggregation to 375 m.\\
Land cover & Copernicus global land-cover 2019, 100 m & Modal class for categorical layers; valid-pixel averages for continuous percent covers.\\
Canopy height & GLAD GEDI--Landsat 2019 map, about 30 m & Valid-pixel averages; map epoch and model limitations are retained.\\
Fuel classes/structure & U.S. LANDFIRE 2016 Remap, 30 m & Modal class for FBFM40/FCCS; averages for numeric encodings, followed by stated physical-unit conversions.\\
Reference dry loading & LANDFIRE FCCS lookup library & Lookup on the original native 30 m codes \emph{before} valid-pixel averaging; coverage fraction stored separately.\\
Fuel moisture & U.S. gridMET 100-hour and 1000-hour dead-fuel moisture, about 4 km, daily & Nearest spatial resampling; values are percent of dry fuel mass. These are not live/fine fuel moisture or hourly measurements.\\
Biome & RESOLVE Ecoregions 2017 polygons & Static broad-scale biome assignment at cell centers; fixed 14-class one-hot vocabulary for trained inputs.\\
\bottomrule
\end{longtable}
Resampling coarser sources to 375 m does not produce independent 375 m weather
observations. Static source years may precede or follow an event; the data do
not establish event-date measured fuel or vegetation. Some environmental
variables are available to the general loader but are not inputs to the saved
19-channel predictor. In particular, a moisture raster in the bundle is not
evidence of a learned moisture effect in that fitted model.

For a target center $x_i$ and returned native-node centers $\eta_j$, weather
mapping is
\begin{equation}
j(i)=\argmin_j\|x_i-\eta_j\|_2,\qquad E_i(t)=V_{j(i)}(t).
\end{equation}
The saved node-index and node-distance rasters retain the mapping. Selection
uses the metric projection, not naive Euclidean distance in latitude/longitude.
For continuous native rasters, a valid-support average can be expressed as
\begin{equation}
\overline v_i=
\frac{\sum_j \omega_{ij}m_jv_j}{\sum_j\omega_{ij}m_j},
\label{eq:validaverage}
\end{equation}
where $\omega_{ij}\ge0$ are reprojection overlap weights and $m_j$ is native
validity. An empty denominator produces unavailable output. A modal class
is the code with greatest aggregated valid weight, not the average of numeric
class codes.

\subsection{Derived meteorological and terrain features}
If wind speed $v$ is in $\mathrm{m\,s^{-1}}$ and meteorological direction
$\theta$ is the direction \emph{from which} wind blows, clockwise from north,
the saved eastward and northward components are
\begin{equation}
u=-v\sin\theta,\qquad v_N=-v\cos\theta.
\end{equation}
Angles are converted from degrees to radians. The negative signs convert
``from'' into vector direction ``toward.'' For hourly preceding-interval
precipitation $P_k$ in mm, trailing precipitation is
\begin{equation}
P_k^{(m)}=\sum_{j=0}^{m-1}P_{k-j},\qquad m\in\{24,168,720\},
\end{equation}
with all $m$ values required. The dry-hour recurrence resets to zero when
$P_k\ge0.1$ mm and otherwise increments the previous count for finite
precipitation. It is a lower bound when a dry spell reaches the beginning of
downloaded history; missing source precipitation keeps the streak unavailable.

On the 31.25 m intermediate terrain grid, let $z$ be elevation and $g_x,g_s$
its finite-difference gradients in eastward and increasing-row (southward)
directions. The source code computes
\begin{equation}
\alpha=\frac{180}{\pi}\arctan\sqrt{g_x^2+g_s^2},\qquad
\beta=\operatorname{atan2}(-g_x,g_s)\pmod{2\pi}.
\end{equation}
$\alpha$ is slope in degrees; $\beta$ is downslope aspect clockwise from
projected-grid north. Near-flat fine pixels with gradient magnitude below
$10^{-5}$ have unavailable aspect. Coarse circular averaging uses
\begin{equation}
e_i=\overline{\sin\beta},\quad n_i=\overline{\cos\beta},\quad
\beta_i=\operatorname{atan2}(e_i,n_i),\quad
\rho_i=\sqrt{e_i^2+n_i^2},
\end{equation}
rather than an arithmetic mean of angles; $\rho_i<10^{-5}$ makes coarse
aspect unavailable. Terrain position index is $z_i-\overline z_{\mathcal N_i}$,
and neighborhood elevation standard deviation is
$\sqrt{\max(\overline{z^2}_{\mathcal N_i}-\overline z_{\mathcal N_i}^2,0)}$,
using available cells of a $9\times9$ neighborhood, nominal width 3375 m.
Vegetated cover is the sum of tree, shrub, grass, and crop percentages with
all four required. Its $3\times3$ neighborhood mean is a coverage proxy, not
a measurement of subpixel fuel connectivity.

\subsection{Temporal support and the completed-day convention}
\label{sec:timesupport}
Snapshot clocks and dynamic feature supports are distinct. Instantaneous
hourly fields are tagged $[t,t]$. Precipitation, rain, radiation, and
evapotranspiration fields summarize the preceding hour $[t-3600,t]$.
Trailing rainfall support begins 24 hours, seven days, or 30 days before $t$.
The loader stores the support start and end as Unix seconds for each channel.
These metadata describe source support, not verified publication availability.

Daily gridMET labels correspond to nominal midnight-to-midnight Mountain
Standard Time, with a fixed 07:00 UTC boundary. Let $s_d$ be the UTC start
of daily source interval $d$. The default \path{completed_day} policy chooses
the latest $d$ with
\begin{equation}
s_d+86400\le t<s_d+2\cdot86400.
\end{equation}
No missing date is extrapolated, and no selection is extended beyond the
downloaded series. The optional \path{containing_day} policy instead chooses
the interval with $s_d\le t<s_d+86400$; a full-day value then includes times
after a within-day cutoff and is suitable only under an explicit retrospective
interpretation. Repeating a selected daily or static value over a sequence
does not turn it into a newly observed hourly quantity.

\subsection{Categorical encoding and independent validity masks}
For categorical source $Z$, freeze ascending distinct class codes
$\mathcal C_Z=(c_1,\ldots,c_J)$. At a known source value $z_i$, indicator
channel $j$ is
\begin{equation}
X_{i,j}=\Ind\{z_i=c_j\}.
\end{equation}
Known membership implies exactly one active indicator, with every indicator's
validity mask true. Missing, absent, unknown-category, or padded membership
gives all-zero indicators and false masks for the entire group. A valid
numeric class code zero is preserved when it belongs to the vocabulary;
numeric zeros are not automatically missing.

Bundle-wide discovery before event/time filtering prevents a changing schema
between subsets. The full recorded loader schema expands 72 source fields
into 741 channels: 65 continuous channels and 676 indicators. This is a
schema width, not a statement that all channels are available for every event.
The fitted predictor freezes a different subset: five continuous channels
and 14 biome indicators, for 19 inputs. The checkpoint saves their order,
categories, groups, and normalization. Re-discovering a smaller vocabulary
at inference would change parameter meaning and is rejected.

The quality codes are 0 valid, 1 out of range, 2 source missing, 3 channel
absent, 4 spatial padding, 5 temporal padding, and 6 unknown category.
Flagged raw observations remain in source records but have false loader
masks. A finite fill value is a storage convenience; its number carries no
missingness semantics. The general loader returns environmental values and
masks with shape $[L,C,H,W]$, current labels and masks with shape $[L,1,H,W]$,
and a collated batch adds a leading dimension $N_b$. Coordinates have shape
$[2,H,W]$, ordered projected $x$ then $y$. The \path{pixel_valid} mask
describes padding, and \path{time_valid} describes padded times; neither is
a substitute for environmental or arrival coverage.

\subsection{Geographic bitmap and array exports}
The image bundle contains 1,380 BMP files with geographic sidecars. Categorical
meaning is in palette indices, not displayed RGB brightness. Exact continuous
values are retained separately in \path{.values.npy}; quantized image colors
must not be used to recover scientific precision. A world file lists
$A,D,B,E,C,F$ with upper-left \emph{pixel-center} coordinates, giving
\begin{equation}
x=A\,c+B\,r+C,\qquad y=D\,c+E\,r+F.
\end{equation}
This differs from the pixel-corner affine transform, whose center evaluation
requires the half-pixel shift. CRS, coordinates, timestamps, range, and mask
metadata must travel with a numeric export.


\section{Step 4: exact next-hour transition targets and the statistical population}
\label{sec:probability-targets}

The probability experiments operate on the reconstructed burned-area bundle.
Consequently, their labels describe transitions in a retrospective arrival-time
reconstruction. They are not independently observed hourly ground-truth
measurements. The code implementing this section is
\path{wildfire_burned_area_hourly/probability_model/prepare_training.py},
\path{wildfire_burned_area_hourly/probability_model/filter_contact.py}, and
\path{wildfire_burned_area_hourly/probability_model/cache.py}.\cite{codeprob}

\subsection{Define each example before defining a model}

Let $e$ denote a fire event, $i=(r,c)$ a raster-cell index, and $t$ an
hourly UTC cutoff. In the global notation of Section~\ref{sec:notation},
$A_e(i)=\widehat a_i$ and $a_e(t)=h_e(t)$ below. Let $A_e(i)$ be the reconstructed arrival time at cell $i$,
measured in hours from that event's recorded origin $t_{0,e}$, and let
$a_e(t)=(t-t_{0,e})/(3600\ \mathrm{s})$ be the elapsed time at cutoff $t$.
At cells with valid reconstructed arrival data, define the cumulative state
\begin{equation}
 B_{e,t}(i)=\Ind\{A_e(i)\leq a_e(t)\}\in\{0,1\}.
 \label{eq:prob-cumulative-state}
\end{equation}
This state remains one after arrival; it does not identify an actively flaming
cell. A currently unburned cell has $B_{e,t}(i)=0$. Its exact next-hour target is
\begin{align}
 Y_{e,t}(i)
 &=\Ind\{B_{e,t}(i)=0,\ B_{e,t+3600\mathrm{s}}(i)=1\}\nonumber\\
 &=\Ind\{a_e(t)<A_e(i)\leq a_e(t)+1\}.
 \label{eq:prob-next-hour-target}
\end{align}
The second expression applies only where arrival time is valid. Thus a target
is positive precisely when arrival lies in the half-open interval $(t,t+1
\text{ hour}]$, interpreted through the event origin. A cell burned exactly
at $t$ is already burned and is excluded from the at-risk training population;
a cell burned exactly at $t+1$ hour gives a positive target.

The preparation code enumerates actual hourly timestamps from the catalog
start through end. It uses \texttt{searchsorted(..., side="left")} to find the
first hourly state satisfying $A_e(i)\leq a_e(t)$. It checks that the current
snapshot is exactly at the recorded cutoff, the next snapshot is exactly
3,600 seconds later and inside the event interval, and every dynamic input
has an interval-end timestamp no later than the current cutoff. Fractional
event origins are retained in the elapsed-time calculation rather than
rounded to an integer hour.

\subsection{Spatial support, masks, and contact}

For an odd side length $K\geq3$, put $R_K=(K-1)/2$ and define the offset set
\begin{equation}
 \mathcal D_K=\{-R_K,\ldots,R_K\}^{2},\qquad
 \mathcal D_K^*=\mathcal D_K\setminus\{(0,0)\}.
\end{equation}
Offsets are enumerated in row order. The source associated with target $i$
and offset $\delta$ is $i+\delta$. This is the PyTorch cross-correlation
convention; no reversal of the offsets is made. Define the number of currently
burned neighboring cells by
\begin{equation}
 n_{e,t,K}(i)=\sum_{\delta\in\mathcal D_K^*}B_{e,t}(i+\delta),
 \qquad C_{e,t,K}(i)=\Ind\{n_{e,t,K}(i)>0\}.
 \label{eq:prob-contact}
\end{equation}
Here $C=1$ means contact with a previously burned cell within the support; it
does not imply physical flame contact or an observed heat-transfer pathway.

Let $M^B_{e,t}(j)$ indicate that burned state is known at source cell $j$.
Let $X_{e,t,c}(j)$ be environmental channel $c$ at that cutoff, and let
$M^X_{e,t,c}(j)$ indicate that its value is valid and finite. Before contact
filtering, a training target must satisfy all of the following conditions:
\begin{enumerate}
 \item Its current and exact next-hour arrival-derived states are known, its
 current state is zero, and the next hour is inside the catalog interval.
 \item Every burned-state source in $i+\mathcal D_K$ is known.
 \item Every selected environmental channel at every position in
 $i+\mathcal D_K$ is known and finite.
 \item The entire square support lies inside the raster. The saved fitting
 contract uses unknown exterior boundaries and complete interior patches.
\end{enumerate}
Let $\mathcal U_K$ denote the finite population of such cell-hours. The
maintained contact-only experiments use the restricted population
\begin{equation}
 \mathcal U_K^+=\{(e,t,i)\in\mathcal U_K:C_{e,t,K}(i)=1\}.
 \label{eq:prob-population}
\end{equation}
The empirical probability target is therefore conditional on complete
reconstruction context and, in the maintained fits, on current burned-neighbor
contact. Moreover, the arrival bundle supplies a final-perimeter-constrained
retrospective support. Evaluation on this population does not estimate the
probability of ignition anywhere in a surrounding landscape.

\subsection{Why the contact restriction changes the task}

Both architectures below sum coefficients multiplied by the current burned
state, with no intercept after that sum. Therefore
\begin{equation}
 n_{e,t,K}(i)=0\quad\Longrightarrow\quad
 \lambda_{e,t}(i)=0,\quad q_{e,t}(i)=0,\quad s_{e,t}(i)=0.
 \label{eq:prob-no-contact-zero}
\end{equation}
This is an algebraic property, independent of training. In the earlier
five-channel unrestricted $3\times3$ experiment, 5,549 of 20,345 positive
test targets had no current burned neighbor: $27.2745\%$ of positive targets.
All of those targets necessarily received probability zero, so the exact
Bernoulli negative log likelihood was infinite. The source for these counts is
\path{probability_model/training_runs/20261007_1h_all44/fitted/metrics.json}
relative to the burned-area package directory. Filtering to
$\mathcal U_K^+$ removes this structurally unsupported part of the task.
It does not repair the model's ability to predict disconnected arrivals; it
defines a narrower task and must be reported as such.

\section{Sample preparation, inverse inclusion weights, and fixed splits}
\label{sec:probability-preparation}

\subsection{Retain positives and sample negative cell-hours}

Preparation first identifies candidate currently unburned cell-hours with
complete arrival-derived burned context. It retains every positive candidate.
Within each fire, negative candidates are divided into two strata: current
contact and current no-contact. For a stratum $h$ with $N_h$ negative
cell-hours, the code draws
\begin{equation}
 m_h=\min(N_h,L_h)
\end{equation}
uniformly without replacement. The recorded limits are $L_h=5{,}000$ for
contact negatives and $L_h=2{,}000$ for no-contact negatives, with seed 7.
Each selected negative has inclusion probability $\pi_u=m_h/N_h$ and
weight $w_u=1/\pi_u=N_h/m_h$. Every positive has $\pi_u=1$ and $w_u=1$.
There is no draw when $N_h=0$. The implementation samples integer time
segments per cell using cumulative segment lengths; this is equivalent to
uniform sampling from the stratum's disjoint cell-hour identities.

Environmental validity is checked after candidate sampling. Invalid selected
examples are removed without replacement draws and without changing their
weights. Contact filtering likewise retains the original weights exactly.
For a fixed finite valid target population $\mathcal V$ and any fixed
quantity $f_u$, let $I_u$ indicate initial sampling. Then
\begin{equation}
 \E_{\rm sampling}\!\left[
   \sum_{u\in\mathcal V}I_uw_uf_u
 \right]
 =\sum_{u\in\mathcal V}f_u,
 \label{eq:prob-ht-total}
\end{equation}
because $\E[I_u]=\pi_u$ and $w_u\pi_u=1$. Filtering a deterministic valid
subset preserves this identity for totals; a replacement draw would change
the inclusion design and require a new calculation.

Most reported means use the ratio
\begin{equation}
 \widehat\mu_f=\frac{\sum_{u\in\mathcal S}w_uf_u}
                    {\sum_{u\in\mathcal S}w_u},
 \label{eq:prob-ratio-mean}
\end{equation}
where $\mathcal S$ is the retained sampled population. This is an
inverse-inclusion-weighted ratio estimate, not an exhaustive full-grid mean.
Equation~\eqref{eq:prob-ht-total} makes the numerator and denominator unbiased
estimators of their respective totals for fixed $f$; it does not make their
ratio exactly unbiased. It also does not establish independence between
neighboring cell-hours or eliminate the effect of model fitting on the same
training sample. The weighted rank and reliability diagnostics below are
also sample-based descriptive estimates, with no reported confidence intervals.

\subsection{Whole-season partitions}

The split function is fixed by country and year before fitting:
\begin{center}
\begin{tabular}{lll}
\hline
Partition & Country-year seasons & Catalog events\\
\hline
Training & Korea 2022; United States 2020 and 2021 & 27\\
Validation & Korea 2023; United States 2022 & 7\\
Test & Korea 2025; United States 2024 & 10\\
\hline
\end{tabular}
\end{center}
All cell-hours of a fire remain in its assigned partition. One model per
kernel pools the two countries; the maintained runs do not fit separate
country-specific models. A catalog fire can contribute no eligible targets
after complete-context and contact restrictions, so catalog event counts and
contributing event counts differ.

For the 19-channel maintained native $3\times3$ cache, the saved counts are:
\begin{center}
\begin{tabular}{lrrrr}
\hline
Partition & Contributing fires & Samples & Positives & $\sum w$\\
\hline
Training & 25 & 205,841 & 99,850 & 3,867,244\\
Validation & 5 & 28,318 & 12,936 & 745,565\\
Test & 10 & 53,292 & 14,796 & 315,920\\
\hline
\end{tabular}
\end{center}
The test sample positive fraction is $14{,}796/53{,}292=0.277640$, whereas
the weighted test prevalence is $14{,}796/315{,}920=0.0468346$. The large
difference follows from retaining all positive targets while sampling
negatives. Treating the sample fraction as the target-population prevalence
would be incorrect. Counts are from
\path{probability_model/training_runs/20261007_1h_biome_onehot_pooled/prepared/metadata.json}.

\section{Input encoding and training-only normalization}
\label{sec:probability-features}

\subsection{The fitted probability models use nineteen channels}

The probability-model input schema selects five continuous source fields:
temperature at 2 m, eastward wind at 10 m, northward wind at 10 m, terrain
slope in degrees, and tree-cover percentage. It also selects
\texttt{biome.biome\_class}. The frozen biome vocabulary contains class
codes $1,\ldots,14$, expanded into fourteen distinct indicator channels:
\begin{equation}
 X_{e,t,5+\ell}(j)=\Ind\{\operatorname{biome}_e(j)=\ell\},
 \qquad \ell=1,\ldots,14.
\end{equation}
For a valid categorical cell, every indicator is exactly zero or one and
their sum is exactly one. A missing biome is invalid; it is not represented
by an all-zero known vector. Vocabulary and ordered channel names are frozen
from the source legend and carried in cache metadata and the checkpoint.
The integers assigned to biome classes are not treated as continuous ordered
measurements. Thus $C=5+14=19$ for these fits, even though the general loader
can expose a broader environmental schema.

\subsection{Weighted moments across training patches}

Index the retained training examples by $u\in\mathcal S_{\mathrm{tr}}$.
Write $X_{u,c,\delta}$ for channel $c$ at offset $\delta$ in example $u$.
Put $W_{\mathrm{tr}}=\sum_{u\in\mathcal S_{\mathrm{tr}}}w_u$. The code calculates
\begin{align}
 \mu_c&=\frac{1}{W_{\mathrm{tr}}K^2}
   \sum_{u\in\mathcal S_{\mathrm{tr}}}w_u
   \sum_{\delta\in\mathcal D_K}X_{u,c,\delta},\\
 v_c&=\frac{1}{W_{\mathrm{tr}}K^2}
   \sum_{u\in\mathcal S_{\mathrm{tr}}}w_u
   \sum_{\delta\in\mathcal D_K}(X_{u,c,\delta}-\mu_c)^2,\\
 \sigma_c&=\begin{cases}\sqrt{v_c},&\sqrt{v_c}\geq10^{-6},\\
                           1,&\sqrt{v_c}<10^{-6},\end{cases}\qquad
 Z_{u,c,\delta}=\frac{X_{u,c,\delta}-\mu_c}{\sigma_c}.
 \label{eq:prob-normalization}
\end{align}
All moments use training examples only. Each target's weight applies equally
to all $K^2$ positions in its patch. These are moments of the weighted patch
distribution, not moments of unique raster cells: overlapping patches can
contain the same physical cell multiple times. The calculation uses float64
accumulation and a second centered pass for variance, then saves float32
means and scales. One-hot channels receive the same normalization rule after
their raw indicator validity is checked.

Validation, test, demonstration, and inference reuse the saved $\mu_c$ and
$\sigma_c$ without refitting. Each kernel has its own native training
population and hence its own saved moments. Center-cropping a $9\times9$
evaluation patch to a smaller support does not change that smaller model's
normalization parameters. This distinction matters in the common comparison.

\section{Environmental kernel generator and the two output formulas}
\label{sec:probability-models}

\subsection{Shared position-dependent coefficient generator}

The generator receives the normalized environmental $K\times K$ context of
each target. With $d_h=16$ hidden channels, define first-layer weights
$U_{h,c,\delta}\in\R$, hidden biases $b_h\in\R$, second-layer weights
$V_{\eta,h}\in\R$, and output biases $d_\eta\in\R$, where
$h=1,\ldots,d_h$, $c=1,\ldots,C$, and $\delta,\eta\in\mathcal D_K$.
At a target $u=(e,t,i)$ the exact center computation is
\begin{align}
 a_{u,h}&=b_h+\sum_{c=1}^{C}\sum_{\delta\in\mathcal D_K}
                         U_{h,c,\delta}Z_{u,c,\delta},\\
 H_{u,h}&=\operatorname{SiLU}(a_{u,h})
          =a_{u,h}\operatorname{sigmoid}(a_{u,h}),\\
 z_{u,\eta}&=d_\eta+\sum_{h=1}^{d_h}V_{\eta,h}H_{u,h}.
 \label{eq:prob-generator}
\end{align}
Here $\operatorname{sigmoid}(a)=1/(1+e^{-a})$. The quantity $d_h$
is hidden width, distinct from the raster's row count $H$.
The first layer is a $K\times K$ convolution and the second a $1\times1$
convolution with $K^2$ outputs. Applying these convolutions across a full grid
produces a different vector $(z_{u,\eta})_\eta$ at every target. The same
learned parameter arrays are shared across targets and events. The generator
reads environmental data; the generated coefficients act on the burned-state
map. The burned map is not an additional input channel to the generator.
Leading time and batch dimensions are independent snapshots; these modules
have no temporal recurrence.

Let $m_\eta=\Ind\{\eta\neq(0,0)\}$. Both models fix their center coefficient
to zero by multiplication by $m_\eta$. The output biases $d_\eta$ are biases
of the coefficient generator; they are not an additive intercept after the
burned-source sum. No normalization makes neighboring coefficients sum to one.

The allocated trainable parameter count is
\begin{equation}
 P_K=d_hCK^2+d_h+d_hK^2+K^2=K^2\{d_h(C+1)+1\}+d_h.
\end{equation}
For $C=19$ and $d_h=16$, $P_K$ is 2,905, 8,041, 15,745, and 26,017 for
$K=3,5,7,9$, respectively. The allocated center-output parameters are
mathematically inactive because of the fixed mask. Holding $d_h$ constant
therefore does not hold model capacity constant across kernel sizes.

\subsection{Original nonnegative rate and exponential-link model}

The original architecture maps logits to nonnegative hourly contributions:
\begin{align}
 k_{u,\eta}&=m_\eta\operatorname{softplus}(z_{u,\eta})
            =m_\eta\log(1+\exp z_{u,\eta}),\\
 \lambda_u&=\sum_{\eta\in\mathcal D_K}k_{u,\eta}B_{e,t}(i+\eta),\\
 q_u(\Delta)&=(1-B_{e,t}(i))\{1-\exp(-\Delta\lambda_u)\},
 \qquad \Delta\geq0\text{ in hours}.
 \label{eq:prob-original}
\end{align}
The neural formulas use numerical rate values expressed in $\mathrm{h^{-1}}$
and numerical elapsed durations expressed in hours. Equivalently, the
physical rate corresponding to a softplus output is that output times
$1\ \mathrm{h^{-1}}$. Each $k_{u,\eta}$ has units of inverse hours per burned neighboring cell,
so $\lambda_u$ is an inverse-hour rate and $\Delta\lambda_u$ is dimensionless.
All fitting examples have $B_{e,t}(i)=0$ and $\Delta=1$. The implementation
evaluates $1-e^{-x}$ as $-\operatorname{expm1}(-x)$ for numerical accuracy near
zero. Its training output is $q_u(1)$, not $\lambda_u$ and not a thresholded
classification.

The exponential formula has a constant-rate survival interpretation. If
$S_u(v)$ is the conditional probability that an initially unburned target
has not transitioned by elapsed time $v$, and one assumes
\begin{equation}
 \frac{dS_u(v)}{dv}=-\lambda_uS_u(v),\qquad S_u(0)=1,
\end{equation}
then separation of variables gives $S_u(v)=e^{-v\lambda_u}$ and transition
probability $1-S_u(\Delta)=q_u(\Delta)$. This is a modeling assumption used
to motivate the link, not a fact established by the reconstruction. During
one call, the current burned map and environmental kernel are held fixed;
the code does not let newly burned cells become new sources within that call.
Although the low-level formula accepts other nonnegative durations, the saved
experiments fit and validate only the one-hour endpoint. Good behavior at a
different duration is not demonstrated by these experiments.

At fixed environmental context, with the target held currently unburned,
$\lambda_u$ and $q_u$ cannot decrease when a known off-center source changes
from unburned to burned, because all contributions are nonnegative. If every source is unburned, both values are exactly zero.
These constraints express the chosen architecture and do not prove a causal
effect of any individual environmental channel.

\subsection{Signed direct-convolution model}

The separate direct experiment removes the coefficient nonnegativity
constraint and the exponential output link:
\begin{align}
 w^{\rm ker}_{u,\eta}&=m_\eta z_{u,\eta}\in\R,\\
 s_u&=(1-B_{e,t}(i))\sum_{\eta\in\mathcal D_K}
                  w^{\rm ker}_{u,\eta}B_{e,t}(i+\eta),\\
 p_u&=\operatorname{clip}_{[0,1]}(s_u)
        =\min\{1,\max\{0,s_u\}\}.
 \label{eq:prob-direct}
\end{align}
For $u=(e,t,i)$, the local clipped output $p_u$ is the global next-hour
new-burn probability $q_i(t)$. It is distinct from the active-burning
marginal $p_i(v)$ in the heat calculation.
Here $w^{\rm ker}$ is a dimensionless kernel coefficient, distinct from the
inverse inclusion weight $w_u$. The raw score $s_u$ can be negative or exceed
one and has no hazard-rate interpretation. The fitted signed model trains
against $s_u$ with an identity link. It applies no softplus, sigmoid,
exponential conversion, or clipping in the training objective. The separate
probability field $p_u$ is used for checkpoint selection and all reported
probability metrics. The direct model rejects durations other than exactly
one hour. Adding a burned source can decrease $s_u$ when its signed
coefficient is negative, so monotonicity in burned-source count is not a
property of this model.

\subsection{Availability is an output of inference}

Numerical zeros used internally to clean invalid inputs are accompanied by
availability masks. They do not turn missing observations into known
unburned cells. The original low-level sum can disregard an unknown source
only when its coefficient is exactly zero; otherwise every contributing
source must be known. Its generator requires all selected environmental
channels throughout the $K\times K$ context. The direct model requires the
complete burned-source context, including positions with a zero coefficient.

The boundary mode \texttt{unknown} invalidates required outside-grid
context. Opting explicitly into \texttt{zero} assumes zero outside-grid
features in the coordinates passed to the low-level model and no exterior
burned-source contribution. Saved wrappers normalize first: exterior zero
in normalized feature coordinates corresponds to each channel's raw training
mean, which need not be zero in physical units. The
training caches use complete interior contexts under \texttt{unknown}.
A known already-burned target has zero new-burn probability even when its
rate or environmental context is unavailable. For the original low-level
model, a known target also has zero new-burn probability at zero duration
even when its rate is unavailable. Probability validity and rate validity
are therefore different objects.

The saved contact-only wrappers impose the fitted domain again at prediction:
an unburned target without a known burned neighbor is returned as unavailable
with a false validity mask and a NaN value, rather than accepted as a
validated zero-probability forecast. The algebraic low-level zero in
Equation~\eqref{eq:prob-no-contact-zero} and this fitted-wrapper availability
rule must not be confused.

\section{Initialization, fitting objective, gradients, and model selection}
\label{sec:probability-fitting}

\subsection{Train-only baseline initialization}

For the original family, the trainer first fits a constant nonnegative
per-neighbor rate $r$ on the native training sample:
\begin{equation}
 \widehat r_{\rm exp}\approx
 \arg\min_{0\leq r\leq20}
 \frac{1}{W_{\mathrm{tr}}}
 \sum_{u\in\mathcal S_{\mathrm{tr}}}w_u
  \{1-e^{-r n_{u,K}}-Y_u\}^{2}.
 \label{eq:prob-rate-init}
\end{equation}
A bounded scalar numerical search fits this baseline; the code explicitly
compares candidate endpoints. For larger supports it scales the search
bracket by the smallest occupied positive neighbor count to avoid a saturated
initial search, while retaining the original admissible endpoint candidates.
This numerical procedure is not a proof that every possible nonconvex scalar
objective attains its global minimum at the returned rate. The final-layer
weights are set to zero and every output bias is set to the inverse softplus
of $\max(\widehat r_{\rm exp},10^{-8})$. Consequently, before training every
off-center coefficient equals this common initialized rate.

For the signed family, a constant linear coefficient has a closed-form
training-only least-squares solution. Let $n_u=n_{u,K}$. Differentiating
\begin{equation}
 F(r)=\sum_{u\in\mathcal S_{\mathrm{tr}}}w_u(rn_u-Y_u)^2
\end{equation}
gives $F'(r)=2\sum_uw_un_u(rn_u-Y_u)$. Since
$\sum_uw_un_u^2>0$ on a nonempty contact-only sample, the unique minimizing
coefficient is
\begin{equation}
 \widehat r_{\rm lin}=
 \frac{\sum_uw_un_uY_u}{\sum_uw_un_u^2}.
 \label{eq:prob-linear-init}
\end{equation}
The allowed coefficient is signed, although this particular baseline is
nonnegative because its numerator is nonnegative. No clipping is used to fit
it. The direct model initializes final-layer weights to zero and output
biases to $\widehat r_{\rm lin}$, then masks the center coefficient to zero.

\subsection{Weighted squared error, with two fixed scaling factors}

Write $N_{\mathrm{tr}}=|\mathcal S_{\mathrm{tr}}|$, and define the global mean
training weight and weighted training prevalence by
\begin{equation}
 \overline w_{\rm tr}=\frac{W_{\mathrm{tr}}}{N_{\mathrm{tr}}},\qquad
 \rho_{\mathrm{tr}}=\frac{\sum_{u\in\mathcal S_{\mathrm{tr}}}w_uY_u}{W_{\mathrm{tr}}}.
\end{equation}
The trainers require $0<\rho_{\mathrm{tr}}<1$. On a minibatch $\mathcal B$ of
$b$ examples, the exact mathematical objective before float32 rounding is
\begin{equation}
 L_{\mathcal B}(\theta)=
 \frac{1}{b\rho_{\mathrm{tr}}}
 \sum_{u\in\mathcal B}\frac{w_u}{\overline w_{\rm tr}}
    \{f_\theta(u)-Y_u\}^{2},
 \quad
 f_\theta(u)=
 \begin{cases}q_u(1),&\text{original family},\\
               s_u,&\text{signed direct family}.
 \end{cases}
 \label{eq:prob-minibatch-loss}
\end{equation}
Both scaling factors are fitted once from the full native training partition;
neither is recalculated within a minibatch. In particular, the denominator
does not sum minibatch weights. With uniformly shuffled training identities,
the expectation of a fixed-size randomly selected batch loss, conditional on
the prepared sample and fixed parameters, equals
\begin{equation}
 \frac{1}{\rho_{\mathrm{tr}}W_{\mathrm{tr}}}
 \sum_{u\in\mathcal S_{\mathrm{tr}}}w_u\{f_\theta(u)-Y_u\}^{2}.
\end{equation}
Thus the two positive global scaling factors change gradient magnitudes but
not the minimizers of the unregularized full-sample weighted squared error.
The sampling weights are not replaced by an arbitrary class-balancing loss.
The optimizer separately uses a recorded weight-decay coefficient.

For an ideal population and unrestricted function class, squared loss has
conditional mean $\E[Y\mid\text{model inputs, eligibility}]$ as its minimizer:
if $m=\E[Y\mid X]$, then
\begin{equation}
 \E[(f-Y)^2\mid X]=(f-m)^2+\operatorname{Var}(Y\mid X).
\end{equation}
This explains the probability interpretation of fitting a binary response
with squared error. The finite neural architecture, selected training sample,
regularization, optimizer, and retrospective label definition still constrain
what the fitted output estimates. Neither this identity nor a low Brier score
establishes calibration of a fitted model.

\subsection{Gradients show exactly which quantity is fitted}

For $\alpha_u=2w_u/(b\overline w_{\rm tr}\rho_{\mathrm{tr}})$, a batch example
has derivative $\partial L/\partial f_u=\alpha_u(f_u-Y_u)$. In the original
family, on a currently unburned valid target, the chain rule gives
\begin{align}
 \frac{\partial q_u(\Delta)}{\partial\lambda_u}
   &=\Delta e^{-\Delta\lambda_u},\\
 \frac{\partial\lambda_u}{\partial z_{u,\eta}}
   &=m_\eta B_{e,t}(i+\eta)\operatorname{sigmoid}(z_{u,\eta}),\\
 \frac{\partial L}{\partial z_{u,\eta}}
   &=\alpha_u(q_u-Y_u)\Delta e^{-\Delta\lambda_u}
      m_\eta B_{e,t}(i+\eta)\operatorname{sigmoid}(z_{u,\eta}).
 \label{eq:prob-original-gradient}
\end{align}
The actual fitting duration is $\Delta=1$. Hence the loss acts on the linked
probability while gradients pass through its rate. It is not squared error
between rate and target.

For the signed identity-link family,
\begin{equation}
 \frac{\partial L}{\partial z_{u,\eta}}
 =\alpha_u(s_u-Y_u)m_\eta B_{e,t}(i+\eta).
 \label{eq:prob-direct-gradient}
\end{equation}
A negative raw score on a positive target still gives a nonzero training
gradient when a burned source contributes. If clipping had been placed inside
the objective, its derivative outside $(0,1)$ would instead be zero almost
everywhere. The code intentionally trains the unclipped score. In both models,
an inactive source or the fixed center yields zero derivative through that
coefficient for this example. Backpropagation through the shared hidden layer
uses
\begin{equation}
 \operatorname{SiLU}'(a)=\operatorname{sigmoid}(a)
   +a\operatorname{sigmoid}(a)\{1-\operatorname{sigmoid}(a)\}.
\end{equation}
For example, if $g_{u,\eta}=\partial L/\partial z_{u,\eta}$, then the data-loss
gradients are
\begin{align}
 \frac{\partial L}{\partial V_{\eta,h}}
    &=\sum_{u\in\mathcal B}g_{u,\eta}H_{u,h},&
 \frac{\partial L}{\partial d_\eta}
    &=\sum_{u\in\mathcal B}g_{u,\eta},\\
 \frac{\partial L}{\partial U_{h,c,\delta}}
    &=\sum_{u\in\mathcal B}
       \left(\sum_{\eta}g_{u,\eta}V_{\eta,h}\right)
       \operatorname{SiLU}'(a_{u,h})Z_{u,c,\delta}.
\end{align}
These formulas describe the differentiable complete-patch training path.
They do not differentiate through availability masks or the discrete process
that constructs reconstructed labels.

\subsection{The Adam update used by the trainers}
Let $\theta^{(k-1)}\in\R^{P_K}$ be the flattened parameter vector before
minibatch update $k$, and let $g^{(k)}=\nabla_\theta L_{\mathcal B_k}$ be
the data-loss gradient above. The current trainers use PyTorch 2.5.1
\path{torch.optim.Adam}, with learning rate $\eta=10^{-3}$ and coupled
weight decay $\lambda_{\mathrm{wd}}=10^{-4}$. No optimizer-specific
defaults are overridden for the moment coefficients or denominator constant,
so $\beta_1=0.9$, $\beta_2=0.999$, and $\epsilon_{\mathrm{Adam}}=10^{-8}$.
These defaults were checked in the locally installed pinned optimizer source.

For each parameter coordinate, the mathematical update is
\begin{align}
\widetilde g^{(k)}&=g^{(k)}+\lambda_{\mathrm{wd}}\theta^{(k-1)},\\
m^{(k)}&=\beta_1m^{(k-1)}+(1-\beta_1)\widetilde g^{(k)},\\
v^{(k)}&=\beta_2v^{(k-1)}+(1-\beta_2)(\widetilde g^{(k)})^2,\\
\widehat m^{(k)}&=m^{(k)}/(1-\beta_1^k),\qquad
\widehat v^{(k)}=v^{(k)}/(1-\beta_2^k),\\
\theta^{(k)}&=\theta^{(k-1)}-
\eta\frac{\widehat m^{(k)}}{\sqrt{\widehat v^{(k)}}+\epsilon_{\mathrm{Adam}}},
\end{align}
starting from $m^{(0)}=v^{(0)}=0$. Products, powers, square roots, and division
are componentwise. The index $k$ counts minibatch updates, not epochs.
Coupled weight decay adds the derivative of
$\lambda_{\mathrm{wd}}\|\theta\|_2^2/2$ to the data gradient before the
moment calculations; it is not an independent multiplicative shrinkage step.
The first and second moment recurrences and their bias corrections adapt the
step to the accumulated gradient history. Floating-point and backend execution
can give small differences from these real-arithmetic equations.

The trainers clear old gradients for each batch, evaluate the appropriate
linked probability or raw score, reject a nonfinite loss, backpropagate,
reject any nonfinite parameter gradient, and perform the optimizer update.
They do not clip gradients in the maintained fitting loop. Validation and
test calculations do not update the optimizer or preprocessing statistics.


\subsection{Epoch selection precedes test evaluation}

The saved runs use Adam, learning rate $10^{-3}$, weight decay $10^{-4}$,
batch size 2,048, hidden width 16, seed 7, at most 30 epochs, and patience 6,
with float32 model calculations on the recorded MPS device. Each epoch
shuffles training examples, performs the minibatch updates, and computes
native weighted validation Brier. A checkpoint is saved only on a strict
improvement, and training stops after six successive epochs without one.
The original family selects by $q_u$; the direct family selects by
$\operatorname{clip}(s_u,0,1)$. The direct raw squared error used for
optimization and clipped squared error used for selection are distinct
quantities. Clipping cannot increase squared error against a binary response,
since projecting onto $[0,1]$ cannot increase distance to either $0$ or $1$.

The selected checkpoint is reloaded before the trainer first loads the test
cache. Selected epochs for $K=3,5,7,9$ were respectively $2,8,5,3$ in the
original family and $4,1,4,6$ in the signed family. Optimization is finite and
nonconvex; no guarantee of a global neural-network minimum is claimed.

The optimized center-patch computation is algebraically the center of the
public full-spatial forward calculation. The trainers verify this agreement
on up to 32 examples in float32 and independently with a CPU float64 model.
Separate float32 tolerances account for different accumulation orders.
Reloaded checkpoints also carry the ordered channel schema, normalization,
one-hour target, zero-center constraint, split assignments, contact rule, and
prepared-metadata hash. Prediction therefore has an explicit saved contract.

\section{A common cohort for comparing kernel sizes and model families}
\label{sec:probability-common}

\subsection{Native populations are different}

Increasing $K$ changes both the complete-context requirement and the contact
requirement. More cells lose full context near a missing-data boundary, while
more distant burned sources can satisfy contact. Native eligible populations
and their sampling weights consequently differ. Even the training-only
normalization distribution differs. Comparing native metrics as though the
models had been evaluated on the same examples is invalid.

The preserved comparison takes the native $9\times9$ prepared sample and
retains only anchors having current burned contact inside its centered
$3\times3$ support:
\begin{equation}
 \mathcal S_{\rm common}
 =\{u\in\mathcal S_9^+:n_{u,3}>0\}.
 \label{eq:prob-common-cohort}
\end{equation}
The complete $9\times9$ patch guarantees that each smaller centered patch
is complete. Membership depends on current context and current burned state;
it does not use the next-hour target. Every model receives the same event,
cutoff, row, column, label, and original $9\times9$ inverse inclusion weight.
For $K<9$, the model reads only the centered $K\times K$ crop and applies its
own saved normalization. There is no intersection of independently sampled
negative rows and no refitting on the common evaluation cohort.

\begin{center}
\begin{tabular}{lrrrrr}
\hline
Common partition & Fires & Samples & Positives & $\sum w$ & Prevalence\\
\hline
Training & 22 & 100,259 & 64,285 & 2,291,049.6466 & 0.0280592\\
Validation & 3 & 12,143 & 6,579 & 343,561.4514 & 0.0191494\\
Test & 8 & 16,807 & 7,378 & 85,133.3236 & 0.0866641\\
\hline
\end{tabular}
\end{center}
The represented populations are weighted estimates of the complete
$9\times9$, current $3\times3$-contact subpopulation. The common validation
cohort contributes only three fires, all in the United States; no Korean fire
survives its complete $9\times9$ context and centered $3\times3$ contact
requirements. The common training evaluation covers two Korean and twenty
US fires, and common test covers two Korean and six US fires. Kernel-size
selection therefore has direct common-validation evidence only from the
retained US fires, despite fitting pooled-country models. No uncertainty
intervals or independent fire-level resampling were supplied with these results.

\subsection{Kernel selection is a separate validation operation}

Each already fitted kernel keeps the epoch selected on its own native
validation population. The common validation population then selects the
kernel size by minimum weighted Brier, breaking exact ties toward the smaller
kernel. The choice is saved before the common test cache is read. Both
families selected $5\times5$. The shared validation Brier values are:
\begin{center}
\begin{tabular}{lrrrr}
\hline
Family & $3\times3$ & $5\times5$ & $7\times7$ & $9\times9$\\
\hline
Original & 0.01868646 & 0.01859248 & 0.01869454 & 0.01867230\\
Signed direct, clipped & 0.01869695 & 0.01859025 & 0.01868570 & 0.01865468\\
\hline
\end{tabular}
\end{center}
Kernel geometry, training population, parameter count, and fitted
normalization all vary across supports. This comparison tests the resulting
trained systems on a common evaluation population; it is not an isolated,
equal-capacity ablation of spatial radius.

\section{Metrics, saved numerical results, and failure modes}
\label{sec:probability-metrics}

\subsection{Weighted Brier and exact Bernoulli log loss}

Sampling weights $w_u$ are indexed by examples $u$ and are distinct from
physical dry fuel loading $w_i$ indexed by cells $i$.
For evaluation examples $u\in\mathcal S$ with positive weights, labels
$Y_u\in\{0,1\}$ and probability outputs $p_u\in[0,1]$, put
$W=\sum_uw_u$. The weighted Brier score is
\begin{equation}
 \widehat{\operatorname{BS}}=\frac{1}{W}\sum_uw_u(p_u-Y_u)^2.
\end{equation}
The original model supplies $p_u=q_u(1)$; the direct model supplies the clipped
view from Equation~\eqref{eq:prob-direct}. The direct raw-score mean squared
error is recorded separately and must not be called a probability Brier score
without explaining that scores may lie outside $[0,1]$.

Define the Bernoulli negative log-likelihood contribution without an ambiguous
$0\log0$ expression:
\begin{equation}
 \ell(Y,p)=\begin{cases}-\log p,&Y=1,\\-\log(1-p),&Y=0,\end{cases}
 \qquad -\log0=+\infty.
\end{equation}
The exact weighted mean is $\sum_uw_u\ell(Y_u,p_u)/W$. A single positive
target with $p_u=0$, or negative target with $p_u=1$, makes it infinite.
This does not imply infinite squared error. It is a support failure under
the fitted Bernoulli probabilities.

The code additionally reports a diagnostic finite log loss using
\begin{equation}
 p_u^{(\epsilon)}=\min\{1-\epsilon,\max\{\epsilon,p_u\}\},
 \qquad \epsilon=10^{-7}.
\end{equation}
Its weighted log loss uses $p_u^{(\epsilon)}$ in place of $p_u$. The finite
diagnostic is not the exact NLL and does not change the fitted model or its
unmodified output probabilities.

\subsection{Weighted ranking metrics, with explicit tie treatment}

Let $P=\sum_uw_uY_u$ and $N=\sum_uw_u(1-Y_u)$. Sort distinct scores in
descending order, and let $\tau_j$ be the $j$th distinct threshold. All
examples tied at a threshold enter simultaneously. Define
\begin{align}
 \operatorname{TP}_j&=\sum_{u:p_u\geq\tau_j}w_uY_u,&
 \operatorname{FP}_j&=\sum_{u:p_u\geq\tau_j}w_u(1-Y_u),\\
 R_j&=\operatorname{TP}_j/P,&
 F_j&=\operatorname{FP}_j/N,\\
 \Pi_j&=\operatorname{TP}_j/(\operatorname{TP}_j+\operatorname{FP}_j).
\end{align}
With $R_0=F_0=0$, the implemented average precision and trapezoidal ROC
area are
\begin{equation}
 \widehat{\operatorname{AP}}=\sum_j(R_j-R_{j-1})\Pi_j,\qquad
 \widehat{\operatorname{AUC}}=\sum_j(F_j-F_{j-1})\frac{R_j+R_{j-1}}{2}.
\end{equation}
AP is undefined when there are no weighted positives; ROC AUC additionally
requires weighted negatives. The code returns unavailable values in these
cases. Ranking measures do not test probability calibration, and clipping
can introduce ties that affect ranking.

\subsection{Fixed-bin reliability}

The descriptive reliability calculation fixes bin edges in advance at
\begin{equation}
 0,\ 0.01,\ 0.02,\ 0.05,\ 0.10,\ 0.20,\ 0.50,\ 1.
\end{equation}
Each bin is left closed and right open, except the final bin includes one.
For a nonempty bin $b$, define
\begin{equation}
 W_b=\sum_{u\in b}w_u,\qquad
 \overline p_b=\frac{\sum_{u\in b}w_up_u}{W_b},\qquad
 \overline y_b=\frac{\sum_{u\in b}w_uY_u}{W_b}.
\end{equation}
The recorded weighted expected calibration error is
\begin{equation}
 \widehat{\operatorname{ECE}}=\sum_{b:W_b>0}\frac{W_b}{W}
                   |\overline p_b-\overline y_b|.
\end{equation}
It is a fixed-bin descriptive statistic, sensitive to the stated bins and
population. It provides no confidence interval and no general calibration
guarantee. The saved metric calculations are in
\path{wildfire_burned_area_hourly/probability_model/metrics.py}.

\subsection{Saved common test results}

The following table reproduces the saved common test comparison, rounded for
readability. Every row uses the same 16,807 identities, 7,378 positive labels,
and weights. ``Zero positive'' counts realized positive labels given exact
probability zero.
\begin{center}
\begingroup
\small
\setlength{\tabcolsep}{3pt}
\begin{tabular}{llrrrrrr}
\hline
Family & $K$ & Brier & AP & ROC AUC & Floor LL & Exact NLL & Zero-pos.\\
\hline
Original & 3 & 0.067219 & 0.383729 & 0.831045 & 0.242273 & 0.242273 & 0\\
Original & 5 & 0.069587 & 0.333213 & 0.825043 & 0.249690 & 0.249690 & 0\\
Original & 7 & 0.071997 & 0.387380 & 0.838929 & 0.259517 & 0.259517 & 0\\
Original & 9 & 0.068417 & 0.363097 & 0.829479 & 0.251728 & 0.251728 & 0\\
Signed direct & 3 & 0.069701 & 0.432513 & 0.823782 & 0.274618 & $+\infty$ & 164\\
Signed direct & 5 & 0.064982 & 0.471449 & 0.842467 & 0.262619 & $+\infty$ & 215\\
Signed direct & 7 & 0.067309 & 0.429701 & 0.840482 & 0.244700 & $+\infty$ & 61\\
Signed direct & 9 & 0.077888 & 0.289826 & 0.767516 & 0.329989 & $+\infty$ & 290\\
\hline
\end{tabular}
\endgroup
\end{center}
No common test row had a negative label assigned probability one. Thus the
infinite exact NLL of every signed direct model is caused by positive labels
receiving probability zero. In these models, negative raw scores are clipped
to zero; contact conditioning does not prevent this different support failure.

For the selected direct $5\times5$ model, saved common test raw scores range
from $-0.0415605344$ to $0.7262458205$. There are 2,162 negative raw scores,
or $12.8637\%$ of sampled targets and $19.5557\%$ of represented population
weight; no raw score exceeds one. Of the 7,378 positive targets, 215 receive
zero probability. Its raw weighted squared error is $0.0650427772$ and its
clipped Brier is $0.0649824060$. Mean clipped probability is $0.0458195$,
below weighted outcome prevalence $0.0866641$. Its fixed-bin ECE is
$0.0408446$, compared with $0.0379574$ for the original $5\times5$ model.
These results do not support describing the direct model as calibrated.

The selected direct $5\times5$ improves common test Brier and ranking relative
to the original $5\times5$, but worsens its floor log loss and has infinite
exact NLL. The direct $3\times3$ and $9\times9$ have worse Brier than their
same-size original counterparts, while $5\times5$ and $7\times7$ improve it.
Therefore the experiment does not establish uniform superiority of the signed
family. Also, although both validation selections chose $5\times5$, the
original $3\times3$ has lower descriptive test Brier than the selected original
$5\times5$; the kernel choice is not revised using that test observation.

\subsection{Result files and reproducibility boundaries}

The primary saved evidence is
\path{probability_model/training_runs/20261007_1h_kernel_sweep/common_metrics.json}
and
\path{probability_model/training_runs/20261007_1h_direct_convolution_signed/common_metrics.json}
relative to \path{wildfire_burned_area_hourly/}. Associated
\texttt{common\_predictions} archives retain labels, weights, cell-hour
identities, per-kernel neighbor counts, and probabilities; direct archives
also retain unclipped raw scores. The saved protocols declare kernel sizes,
source/cache identities, training parameters, clipping behavior, and
validation-only selection. The direct evaluation code checks exact equality
of all shared identities, labels, weights, and neighbor counts against the
preceding original common-cohort archives.

The saved common test public-wrapper versus center-computation checks have
maximum absolute discrepancies no greater than $1.79\times10^{-7}$ for the
original family and $5.97\times10^{-8}$ for the signed family, below their
respective recorded float32 tolerances. These are computational agreement
checks. They do not validate the reconstructed label against an independent
observation or quantify generalization uncertainty. The experiments use one
fixed sampling seed and one fitted run per configuration; the reported
comparisons are retrospective, population-conditioned, and descriptive.

\section{Step 5: ecological fuel loading and its coverage contract}
\label{sec:ecological-fuel}

\subsection{Two different reference catalogs have two different roles}

The radiation extension uses FCCS dry loading to prescribe a conditional heat
source, whereas the residence extension uses the Scott--Burgan FBFM40 catalog
to prescribe a passing-front duration. These are separate inputs. A code in one
catalog is not a code in the other catalog. A biome class is neither a dry mass
nor a particle diameter. No conversion from a biome label, canopy height,
relative humidity, or soil moisture to dry fuel mass is implemented by
\texttt{EcologicalFuelMaps}. The implemented contracts are in
\path{expected_heat_flux/ecology.py} and
\path{expected_heat_flux/residence.py}.\cite{coderesidence}

LANDFIRE assigns vegetation to reference FCCS fuelbeds. The reference loading
table describes fuelbed attributes, rather than a field measurement of the
fuel available on the date of a particular fire.\cite{landfirefccs}
FCCS represents canopy, shrub, herbaceous, woody, litter--lichen--moss, and
ground-fuel strata; a reference fuelbed may need customization for an actual
site.\cite{prichard2013} The package reads the previously prepared total-loading
raster. It does not implement all FCCS fire-behavior, consumption, or emissions
calculations.

\subsection{Objects, units, and conversion to SI}

Let $i$ denote one cell of a target raster with shape $[H,W]$, where $H$ and $W$
are positive integer numbers of rows and columns. Let
$w_i^{\mathrm{mapped}}$ be the mean reference dry loading over the mapped
portion of cell $i$, in $\mathrm{kg\,m^{-2}}$ of that portion, and let
$c_i\in[0,1]$ be the fraction of the cell described by the mapped loading.
An unreadable loading or coverage is represented as missing, not as zero.
The distinction between a mapped-part mean and a whole-cell mean is essential:
if $c_i=1/2$, then $w_i^{\mathrm{mapped}}$ is the mean over half of the cell.
Multiplication by $c_i$ yields its contribution per unit \emph{whole-cell} area.

The raster loading is interpreted as U.S. short tons per international acre.
One U.S. short ton is $2000$ international pounds, one international pound is
exactly $0.45359237\,\mathrm{kg}$, one international acre is
$43560$ square international feet, and one international foot is exactly
$0.3048\,\mathrm{m}$. Hence the implemented conversion factor is
\begin{align}
 C_w
 &=\frac{2000\times0.45359237}
          {43560\times(0.3048)^2}
   \frac{\mathrm{kg\,m^{-2}}}{\mathrm{short\ ton\,acre^{-1}}}
   \nonumber\\
 &=0.2241702312388912
   \frac{\mathrm{kg\,m^{-2}}}{\mathrm{short\ ton\,acre^{-1}}},
 \qquad
 w_i^{\mathrm{mapped}}=C_w L_i,
 \label{eq:fccs-si}
\end{align}
where $L_i$ is the loading quantity read from the raster, expressed in
U.S. short tons per international acre; the stored scalar is its numerical
magnitude in those units. This is a unit
conversion; it supplies no information about the event's actual consumed mass.

\subsection{Alignment, physical validity, and complete coverage}

The loading and coverage rasters must each have one static band. Their shapes,
coordinate reference systems (CRSs), and six affine-transform coefficients must
agree exactly. Their CRS must be projected with a linear unit equal to one
metre. If a prediction grid or an expected shape is supplied, the reader checks
that alignment as well. Shape equality by itself is insufficient: a shifted
grid would associate the loading with a different location.

The reader replaces a negative or nonfinite loading by missing data. It also
replaces a nonfinite coverage or a coverage outside $[0,1]$ by missing data.
The numerical tolerance for declaring coverage complete is a caller-selected
$\varepsilon_c\in[0,1)$, whose default is $10^{-6}$. The reader's complete-input
indicator is
\begin{equation}
 V_i^{\mathrm{fuel}}
 =\Ind\{w_i^{\mathrm{mapped}}\text{ is finite}\}
  \Ind\{c_i\text{ is finite}\}
  \Ind\{c_i\geq1-\varepsilon_c\}.
 \label{eq:fccs-complete}
\end{equation}
Here $\Ind\{\mathcal A\}$ equals one if statement $\mathcal A$ is true and zero
otherwise. A finite, nonnegative loading of zero is legitimate when it was
actually supplied with valid coverage; missing loading is a different state.
The accepted tolerance concerns numerical coverage. It does not quantify
uncertainty in fuel mass. Setting $\varepsilon_c=0$ requires exact coverage
$c_i=1$.

Without an explicit fallback loading, the conversion to combustion parameters
requires every cell to satisfy \eqref{eq:fccs-complete}; otherwise it fails. If a
coverage just below one is accepted within tolerance and no fallback is given,
the implementation uses $w_i^{\mathrm{mapped}}$ directly as the whole-cell
loading. This is the precise numerical tolerance convention used by the code.

\subsection{Explicit completion of the unmapped fraction}

Let $w_i^{\mathrm{fallback}}\geq0$ be a finite loading supplied by the caller for
the \emph{unmapped} portion, in $\mathrm{kg\,m^{-2}}$ of that portion. A scalar
fallback applies to every cell; an array must broadcast to the existing spatial
shape $[H,W]$, without introducing a batch or time dimension. When a fallback
is supplied, the whole-cell dry loading is
\begin{equation}
 w_i=
 \begin{cases}
 c_i w_i^{\mathrm{mapped}}
 +(1-c_i)w_i^{\mathrm{fallback}},
 &w_i^{\mathrm{mapped}}\text{ and }c_i\text{ are readable},\\[2pt]
 w_i^{\mathrm{fallback}},
 &\text{otherwise}.
 \end{cases}
 \label{eq:fccs-fallback}
\end{equation}
For a cell of area $|S_i|$ in $\mathrm{m^2}$, the first case follows from adding
the dry masses $|S_i| c_i w_i^{\mathrm{mapped}}$ and
$|S_i|(1-c_i)w_i^{\mathrm{fallback}}$, then dividing by $|S_i|$. If $c_i=1$, the
fallback contributes nothing. If $c_i=0$, the fallback supplies the whole-cell
loading. If either mapped input is unreadable, the implementation uses the
fallback for the entire cell; it cannot recover a known mapped contribution
from unreadable data.

For example, if $c_i=0.4$, $w_i^{\mathrm{mapped}}=2$, and
$w_i^{\mathrm{fallback}}=10$ in $\mathrm{kg\,m^{-2}}$, then
$w_i=0.4(2)+0.6(10)=6.8\,\mathrm{kg\,m^{-2}}$. Replacing the whole cell by the
fallback would discard the mapped information; using the mapped mean alone
would leave the unmapped fraction unspecified.

Under the \emph{reference-input assumptions}, and assuming only that the
unmapped loading is nonnegative, a mathematical lower bound is
\begin{equation}
 w_i^{\mathrm{lower}}=
 \begin{cases}
 c_i w_i^{\mathrm{mapped}},&\text{both mapped inputs are readable},\\
 0,&\text{otherwise}.
 \end{cases}
 \label{eq:fccs-lower}
\end{equation}
Indeed, the unknown nonnegative mass contributes a nonnegative summand in
\eqref{eq:fccs-fallback}. This is not a lower bound on true field fuel mass,
because the reference loading itself can be inaccurate. The zero in the second
case is the trivial nonnegative bound, not an observed zero loading.

\subsection{Dry fuel, moisture, and prescribed combustion parameters}

For a fuel sample with dry mass $m_{\mathrm{dry}}>0$ and water mass
$m_{\mathrm{water}}\geq0$, dry-basis moisture would be
$M=m_{\mathrm{water}}/m_{\mathrm{dry}}$ (dimensionless), or $100M$ percent. That
physical definition is distinct from relative humidity, soil-water content,
and a categorical biome label. The ecological reader neither derives
dead-fuel moisture nor estimates live-fuel moisture. Weather and fuel-moisture
fields can be predictor covariates, but the heat extension does not
automatically turn them into a moisture correction of dry loading, combustion
heat, consumed fraction, or residence time.

After specifying $w_i$, the caller must also specify a positive effective heat
of combustion $H_{c,i}$ in $\mathrm{J\,kg^{-1}}$, a positive combustion-phase
duration $T_i$ in seconds, and a consumed dry-mass fraction $f_i\in[0,1]$.
The conditional uniform heat-release-rate density is
\begin{equation}
 Q_i=\frac{w_i f_i H_{c,i}}{T_i}
 \quad[\mathrm{W\,m^{-2}}].
 \label{eq:ecology-conditional-hrr}
\end{equation}
The effective heat is supplied, rather than inferred as a moisture-corrected
gross calorific value. If $T_i$ is a main-front flaming duration, $f_i$ must
refer to dry fuel consumed in that modeled phase; it must not silently mean
all fuel eventually consumed by residual flaming and smouldering. A separate
radiative fraction is required by the radiation operator.

The completed integrated FishCreek demonstration prescribes
$H_c=18\times10^6\,\mathrm{J\,kg^{-1}}$, $T=300\,\mathrm{s}$,
$f=0.10$, an unmapped-loading fallback of
$10\,\mathrm{kg\,m^{-2}}$, and radiative fraction $0.30$.
These are illustrative scenario parameters in its saved summary. They are not
estimated from observations of FishCreek combustion. The $300\,\mathrm{s}$
duration in that demonstration is separate from the FBFM40 passing-front
durations used by the active-flame demonstration described below.

\section{Fuel-dependent surface-flame residence scenarios}
\label{sec:residence}

\subsection{What duration is being assigned}

In this section, duration is denoted $\tau_i^{\mathrm{fl}}$ to distinguish
it from a separately supplied heat-emission duration $T_i$.
At a representative grid location $i$, let $\tau_i^{\mathrm{fl}}>0$ be a deterministic, finite
duration in seconds. Anderson's passing-front approximation and its Andrews
formulation describe the time for a surface-fire flaming front to pass one
fixed fuel location. The formulation excludes the indefinite trailing edge and
assumes no effect of fuel-bed packing.\cite{anderson1969,andrews2018}
It is distinct from the time since ignition, the time to traverse a coarse
cell, the interval during which some part of that cell contains any flame,
residual-flame extinction, crown-fire duration, or smouldering duration.

The object \texttt{FlameResidenceTime} stores either one scalar duration or a
complete matrix of shape $[H,W]$, plus immutable provenance. A scalar declares
the same scenario duration at all locations to which it is applied. A matrix
must align with the consumer's spatial grid. Zero, negative, nonfinite, masked,
boolean, string, and complex durations are rejected; a missing duration is not
assigned zero. The class copies the input and freezes its duration and
provenance. Direct construction requires a provenance mapping.

Assigning a point-scale $\tau_i^{\mathrm{fl}}$ to a $375\,\mathrm{m}$ grid cell is an additional
spatial modeling assumption. It assigns a duration to a representative
location; it does not resolve the arrival progression or flame geometry within
the footprint. A catalog assignment supplies a reference-fuel scenario, not an
event-specific measurement of duration.

The heat-emission duration $T_i$ in
\eqref{eq:ecology-conditional-hrr} and the flame residence
$\tau_i^{\mathrm{fl}}$ are separate objects. The helper
\texttt{CombustionParameters.from\_flame\_residence} explicitly sets
$T_i=\tau_i^{\mathrm{fl}}$ when constructing combustion parameters from the
residence scenario. When \texttt{WildfireHeatModel} is also supplied a
\texttt{residence\_time}, it enforces this equality and rejects a mismatch.
The equality is a chosen bridge between modeled phases, not a conclusion that
all reference fuel is consumed during front passage; the consumed fraction
must describe that same phase. Without this explicit bridge, an independently
supplied heat duration and a flame residence must retain their distinct
meanings.

\subsection{From particle diameter: the Anderson relation}

Let $d_i>0$ be the particle diameter in metres. Define the dimensionless
original-unit numerical values $d_{\mathrm{inch}}=d_i/(0.0254\,\mathrm m)$
and $\tau_{\mathrm{min}}^{\mathrm{fl}}=\tau_i^{\mathrm{fl}}/(60\,\mathrm s)$.
The implemented source law is
$\tau_{\mathrm{min}}^{\mathrm{fl}}=8d_{\mathrm{inch}}$. Converting units gives
\begin{align}
 d_{\mathrm{inch}}&=\frac{d_i}{0.0254\,\mathrm{m}},\\
 \tau_i^{\mathrm{fl}}&=8\frac{\mathrm{min}}{\mathrm{inch}}
       \frac{d_i}{0.0254\,\mathrm{m}/\mathrm{inch}}
       \frac{60\,\mathrm{s}}{\mathrm{min}}
      =C_d d_i,
 \quad
 C_d=18\,897.63779527559\,\mathrm{s\,m^{-1}}.
 \label{eq:residence-diameter}
\end{align}
Thus a particle of diameter one inch, $d_i=0.0254\,\mathrm{m}$, gives
$\tau_i^{\mathrm{fl}}=480\,\mathrm{s}$. A particle of diameter $0.002\,\mathrm{m}$ gives
$\tau_i^{\mathrm{fl}}=37.79527559055118\,\mathrm{s}$. The law approximates a uniform fuel array;
it does not establish that every fuel particle in a heterogeneous cell burns
out by this time. The public Anderson constructor validates positivity and
finiteness but does not impose a calibrated species-specific diameter range.
Allowing a numerical input therefore does not validate its physical use.

\subsection{From characteristic surface-area-to-volume ratio}

Let $\sigma_i>0$ be characteristic fuel-particle surface area divided by
particle volume, in $\mathrm{m^{-1}}$. Its numerical value in
$\mathrm{ft^{-1}}$ is $0.3048$ times its numerical value in
$\mathrm{m^{-1}}$. Andrews's residence relation is
$\tau_{\mathrm{min}}^{\mathrm{fl}}=384/\sigma_{\mathrm{ft^{-1}}}$.\cite{andrews2018}
The implemented SI form is consequently
\begin{equation}
 \tau_i^{\mathrm{fl}}=\frac{C_\sigma}{\sigma_i},\qquad
 C_\sigma=\frac{384\times60}{0.3048}
          \,\mathrm{s\,m^{-1}}
 =75\,590.55118110236\,\mathrm{s\,m^{-1}}.
 \label{eq:residence-sav}
\end{equation}
The unit check is
$(\mathrm{s\,m^{-1}})/(\mathrm{m^{-1}})=\mathrm{s}$.
Using $384/\sigma_i$ with an SI input would omit both conversions.

For a cylinder of diameter $d_i$ and length $\ell>0$, neglecting its end area,
the side area is $\pi d_i\ell$ and the volume is $\pi d_i^2\ell/4$. Hence
$\sigma_i=4/d_i$, so
\begin{equation}
 \frac{C_\sigma}{\sigma_i}
 =\frac{C_\sigma d_i}{4}=C_d d_i,
 \qquad
 \frac{384\times60}{4\times0.3048}
 =\frac{8\times60}{0.0254}.
\end{equation}
The equality is exact because $0.3048=12(0.0254)$. This geometric argument
depends on the long-cylinder approximation. It is not a general conversion
from leaf geometry, tree height, or raster width to a particle diameter.

\subsection{From the FBFM40 reference catalog}

The catalog constructor supports exactly the 40 burnable standard numeric
identifiers shown in Table~\ref{tab:fbfm-groups}. The components of each model
are dead 1-hour, dead 10-hour, dead 100-hour, live herbaceous, and live woody
fuel. These time-lag class names classify fuel components; they are not assigned
flaming durations. The code transcribes the dry loads and particle SAV values
from Scott--Burgan table~7.\cite{scottburgan2005}

\begin{table}[htbp]
\centering
\begin{tabular}{lll}
\hline
Reference group & Supported numeric identifiers & Number\\
\hline
Grass (GR) & $101$--$109$ & $9$\\
Grass--shrub (GS) & $121$--$124$ & $4$\\
Shrub (SH) & $141$--$149$ & $9$\\
Timber understory (TU) & $161$--$165$ & $5$\\
Timber litter (TL) & $181$--$189$ & $9$\\
Slash/blowdown (SB) & $201$--$204$ & $4$\\
\hline
\end{tabular}
\caption{The complete set of burnable FBFM40 codes accepted by the residence
catalog constructor. Numeric gaps between groups are not additional supported
models.}
\label{tab:fbfm-groups}
\end{table}

For one reference fuelbed, index its components by $j$. Let $w_j\geq0$ be the
dry loading in $\mathrm{kg\,m^{-2}}$, $\rho_j>0$ the particle material density
in $\mathrm{kg\,m^{-3}}$, and $\sigma_j>0$ its particle SAV in
$\mathrm{m^{-1}}$.
The particle volume per unit ground area is $w_j/\rho_j$. Its surface area per
unit ground area is consequently
\begin{equation}
 a_j^{\mathrm{surf}}=\frac{w_j\sigma_j}{\rho_j}.
 \label{eq:component-surface}
\end{equation}
The ratio $a_j^{\mathrm{surf}}$ has units of fuel surface area per ground
area, $\mathrm{m^2}/\mathrm{m^2}$, hence is dimensionless; it is distinct
from the footprint set $S_j$. Surface-area weighting means taking the arithmetic mean of the particle SAVs
with weights proportional to these surfaces, not to dry loads alone:
\begin{equation}
 \overline\sigma
 =\frac{\sum_j a_j^{\mathrm{surf}}\sigma_j}{\sum_j a_j^{\mathrm{surf}}}
 =\frac{\sum_j w_j\sigma_j^2/\rho_j}
        {\sum_j w_j\sigma_j/\rho_j}.
 \label{eq:characteristic-sav-general}
\end{equation}
This collapses the live/dead category weighting and within-category surface
weighting described by Andrews into one expression.\cite{andrews2018}
Scott--Burgan uses one common particle density of
$32\,\mathrm{lb\,ft^{-3}}$, dead 10-hour SAV of $109\,\mathrm{ft^{-1}}$, and
dead 100-hour SAV of $30\,\mathrm{ft^{-1}}$.\cite{scottburgan2005}
The common density and any common loading-unit conversion cancel, giving the
implemented formula
\begin{equation}
 \overline\sigma_{\mathrm{ft^{-1}}}
 =\frac{\sum_{j=1}^{5}w_j\sigma_{j,\mathrm{ft^{-1}}}^{\,2}}
        {\sum_{j=1}^{5}w_j\sigma_{j,\mathrm{ft^{-1}}}},
 \qquad
 \overline\sigma_i=\frac{\overline\sigma_{\mathrm{ft^{-1}}}}{0.3048}
 \,\mathrm{m^{-1}},
 \qquad
 \tau_i^{\mathrm{fl}}=\frac{384\times60}{\overline\sigma_{\mathrm{ft^{-1}}}}\,\mathrm{s}.
 \label{eq:fbfm-residence}
\end{equation}
The subscript on a SAV numerical value declares its unit. Entries with zero
loading contribute zero to both sums. In particular, the source-table SAV
placeholder $9999$ in a zero-load component is not physically weighted as
fine fuel.

For GR1, only the dead 1-hour and live-herbaceous entries have positive load:
$w=(0.10,0.30)$ in short tons per acre, with particle SAVs
$(2200,2000)\,\mathrm{ft^{-1}}$. Thus
\begin{align}
 \overline\sigma_{\mathrm{ft^{-1}}}
 &=\frac{0.10(2200)^2+0.30(2000)^2}
         {0.10(2200)+0.30(2000)}
   =\frac{1\,684\,000}{820}
   =2053.658536585366,\\
 \tau^{\mathrm{fl}}&=\frac{23040}{2053.658536585366}\,\mathrm{s}
   =11.219002375296912\,\mathrm{s}.
\end{align}
The published rounded characteristic SAV, $2054\,\mathrm{ft^{-1}}$, is useful
for checking the transcription, but the constructor calculates the weighted
ratio from the component parameters before converting to duration.

The GS4 source has a recorded inconsistency: its component values in table~7
imply approximately $1631\,\mathrm{ft^{-1}}$, consistent with Andrews's
table~8c, while Scott--Burgan's individual GS4 summary prints
$1674\,\mathrm{ft^{-1}}$. The implementation consistently uses the component
table calculation and records the discrepancy in provenance. It does not
silently substitute the inconsistent summary value.

\subsection{Why curing does not change this particular characteristic SAV}

Let a live-herbaceous component have dry loading $w_h$, SAV $\sigma_h$, and
density $\rho_h$. Let $\alpha\in[0,1]$ be the fraction transferred to a new dead
herbaceous component. Scott--Burgan preserves particle SAV during that transfer;
it does not simply combine the transferred material with a dead 1-hour
component of another SAV.\cite{scottburgan2005}
If the transfer also retains the common catalog density, then the contributions
to the denominator of \eqref{eq:characteristic-sav-general} are
\begin{equation}
 \frac{(1-\alpha)w_h\sigma_h}{\rho_h}
 +\frac{\alpha w_h\sigma_h}{\rho_h}
 =\frac{w_h\sigma_h}{\rho_h},
\end{equation}
and the corresponding numerator contributions are
\begin{equation}
 \frac{(1-\alpha)w_h\sigma_h^2}{\rho_h}
 +\frac{\alpha w_h\sigma_h^2}{\rho_h}
 =\frac{w_h\sigma_h^2}{\rho_h}.
\end{equation}
Therefore both sums, $\overline\sigma$, and this SAV-based $\tau^{\mathrm{fl}}$ are invariant to
$\alpha$. The constructor can compute this duration without choosing a live
moisture or curing fraction. This algebraic invariance is specific to this
duration calculation. It does not imply that moisture and curing have no effect
on ignition, spread, heat release, or consumption. The package does not
estimate a curing fraction from moisture as part of the residence constructor.

\subsection{Unsupported, missing, and nonburnable codes}

The input codes must be numeric integer-valued identifiers in a scalar or a
nonempty spatial matrix. A finite noninteger code is rejected even when a
fallback is supplied. Missing, masked, custom, and unsupported integer codes,
including standard nonburnable codes $91,92,93,98,99$, have no residence assigned
by the burnable catalog. If any such location occurs, the constructor fails
unless the caller explicitly supplies a strictly positive, finite
\texttt{fallback\_sav\_m\_inv}. That fallback must be a scalar or match the
code matrix shape exactly.

On a supported cell the catalog SAV is retained. On an unsupported or missing
cell the fallback SAV is substituted and \eqref{eq:residence-sav} is applied.
Overriding a nonburnable class in this way is a caller's combustion scenario;
it is not a duration attributed to the mapped class and is not evidence that
the cell has combustible fuel. Provenance separately records catalog-cell,
fallback-cell, and missing-code counts, finite unsupported codes, and the
fallback value. No fallback is implicit in the public catalog constructor.

The active-flame example \emph{does} explicitly set
$\sigma_{\mathrm{fallback}}=2000\,\mathrm{ft^{-1}}
=6561.679790026246\,\mathrm{m^{-1}}$, giving $\tau^{\mathrm{fl}}=11.52\,\mathrm{s}$ on fallback
locations. This is an example argument, not a universal residence-time default.

\subsection{A distinct laboratory relation for jarrah round wood}

\texttt{from\_eucalypt\_wood\_diameter} implements Burrows's thesis equation~6-1:
\begin{equation}
 \tau_i^{\mathrm{fl}}=0.871\left(\frac{d_i}{1\,\mathrm{mm}}\right)^{1.875}\,\mathrm{s},
 \qquad 1\,\mathrm{mm}\leq d_i\leq16\,\mathrm{mm}.
 \label{eq:burrows-residence}
\end{equation}
The dimensionless ratio in the power makes the unit convention explicit.
This numerical law uses a \emph{millimetre} input, while
\eqref{eq:residence-diameter} uses a metre input. The constructor rejects
extrapolation outside the closed experimental interval.

The original methods describe individual oven-dried jarrah
(\emph{Eucalyptus marginata}) round-wood particles cut to $10\,\mathrm{cm}$,
ignited over a pilot flame, then allowed to burn after the pilot was
extinguished. The measured endpoint is flame-out, distinguished from
cessation of glowing combustion.\cite{burrows1994}
This fuel-specific particle flame-out relation is separate from Anderson's
passing-front approximation. Neither law is selected merely from a biome label.
The thesis includes still-air laboratory context and ambient temperatures of
$26$--$32\,^{\circ}\mathrm C$ in the surrounding experimental description;
the code records these conditions as scenario metadata. Transferring the law
to another species, wet fuel, wind, a whole fuelbed, or a coarse cell requires
additional evidence that the present package does not supply.

The constants $0.871$ and $1.875$ are treated as fixed in the implemented
scenario. Although the thesis prints regression standard errors, the package
does not sample those parameters, propagate their uncertainty, or construct a
flame-duration confidence interval. The exact implemented equation is supported
by the 1994 thesis; the source register distinguishes this from the later
2001 article, whose full equation was not recorded as inspected.

\subsection{Deterministic durations and the limits of uncertainty claims}

There is no implemented residence-time probability distribution, moisture-based
duration fit, quantile sampler, or Monte Carlo propagation of ecological
parameter uncertainty in these modules. A duration fixed by a scenario can
formally be represented by a point mass at $\tau_i^{\mathrm{fl}}$, but that mathematical
representation is not a fitted uncertainty model. The within-hour uniform
arrival assumption used by the forecast bridge concerns arrival time, not
uncertainty in $\tau_i^{\mathrm{fl}}$. A spatial collection of different prescribed durations
is also not a probability distribution of the unknown duration at a given cell.

Likewise, an arrival-time standard deviation alone does not identify an arrival
distribution or its quantiles. A mean and variance can be shared by many
distributions with different tail probabilities. None of the residence
constructors converts a kriging standard-deviation raster into a normal,
lognormal, or calibrated arrival law. Conditional source-strength assumptions
and fuel uncertainty therefore remain outside the computed heat bounds unless
an additional model is supplied explicitly.

\subsection{Saved FishCreek scenario and verification scope}

The active-flame demonstration uses the $2016$ LANDFIRE modal FBFM40 class in
each $375\,\mathrm{m}$ cell for a $2024$ event. The saved grid has
$H=103$, $W=73$, and $HW=7519$ cells. It assigns $7346$ durations from burnable
catalog classes and $173$ durations from the explicit fallback for codes $98$
and $99$. Its assigned durations have minimum
$11.219002375296913\,\mathrm{s}$, median
$18.83071682602384\,\mathrm{s}$, and maximum
$18.83071682602384\,\mathrm{s}$. Those numbers describe the saved reference-fuel
scenario, not measured flame durations.

For the same event, $2642$ saved valid arrival-standard-deviation values have
minimum $32.776554107666016\,\mathrm h$, median
$87.61341857910156\,\mathrm h$, and maximum
$246.66561889648438\,\mathrm h$. These exact raster values show that the
reported arrival uncertainty is much larger than the selected point durations.
The nominal current-flame map at the saved cutoff has an empty active-position
selection. This outcome cannot establish absence of observed flames: the
calculation uses reconstructed point arrivals and does not integrate their
uncertainty or represent residual flames.

Current verification of the residence and ecological-reader contracts passed
$118$ selected tests. They check the original and SI unit equations, cylinder
geometry, all $40$ rounded characteristic SAVs, zero-load placeholders,
immutable residence provenance, unsupported-code fallbacks, Burrows's tested
range, metric raster alignment, incomplete coverage, nonnegative lower bounds,
and the exact coverage-weighted fallback formula. These are mathematical and
software checks. They are not observational calibration of duration, fuel
consumption, active-flame positions, or wildfire heat exposure.

\section{Step 6: activity conversion and expected incident radiant heat}
\label{sec:heat}

This section documents the methods implemented in
\path{expected_heat_flux/model.py}, \path{expected_heat_flux/temporal.py},
\path{expected_heat_flux/pipeline.py}, and \path{expected_heat_flux/flames.py}.
The result is the expected \emph{incident radiant irradiance} at specified
receiver points, in $\mathrm{W\,m^{-2}}$, under an explicit combustion, timing,
and geometry scenario. It is not an estimate of every component of wildfire
heat transfer. In particular, convection, receiver absorption, temperature,
conduction, terrain obstruction, smoke-dependent attenuation, and emission
from outside the supplied source raster are absent from this implementation.
The mathematical calculations can be verified without establishing that the
physical assumptions accurately describe a particular wildfire.\cite{codeheat}

\subsection{Objects, units, and the source-state interpretation}
\label{sec:heat-objects}

Let $\mathcal J=\{0,\ldots,H-1\}\times\{0,\ldots,W-1\}$ be the finite
source-cell index set, where $H,W$ are positive integers. Index $j=(r,c)$
denotes row $r$ and column $c$, with both indices starting at zero. A projected
affine grid in metres has coefficients $(a,b,c_0,d,e,f_0)$ and pixel-corner map
\begin{equation}
 (x,y)=(ac+br+c_0,\;dc+er+f_0).
\end{equation}
The column and row edge vectors are $\mathbf u=(a,d)$ and $\mathbf v=(b,e)$,
respectively. Each cell has horizontal ground area
\begin{equation}
 A_{\rm cell}=|ae-bd|>0 \qquad[\mathrm{m^2}],
\end{equation}
and centre
\begin{equation}
 \mathbf x_j=(a(c+\tfrac12)+b(r+\tfrac12)+c_0,
              d(c+\tfrac12)+e(r+\tfrac12)+f_0).
\end{equation}
This representation supports rotation, shear, rectangular spacing, and either
row direction. Longitude and latitude are angular coordinates and cannot be
used as metre coordinates. The radiation calculation requires a projected
metric grid; $\texttt{AffineGrid}$ itself does not inspect a coordinate-system
identifier. The caller or data adapter must establish that requirement.

Let $\mathcal F_0$ denote the information used at forecast time zero. For
cell $j$, let $Z_j(t)\in\{0,1\}$ indicate active burning at relative time
$t$ seconds, and let
\begin{equation}
 p_j(t)=\mathbb P(Z_j(t)=1\mid\mathcal F_0)
\end{equation}
be its active-burning marginal probability. Throughout the timing bridge,
$t$ is numeric seconds relative to the forecast cutoff. The indicator activates the
model's prescribed uniform source footprint over the entire cell. This is a
spatial idealization: partial active coverage must be represented through the
conditional whole-cell source density, rather than interpreted as a measured
flame filling the raster cell. A cumulative burned-state indicator and a
probability of first burning sometime during the next hour are different
objects from $Z_j(t)$ and $p_j(t)$.

The combustion inputs for cell $j$ are dry fuel loading $w_j\ge0$ in
$\mathrm{kg\,m^{-2}}$, phase-specific consumed fraction $f_j\in[0,1]$,
effective released heat per fuel mass $H_{c,j}>0$ in $\mathrm{J\,kg^{-1}}$, and
constant-emission duration $T_j>0$ in seconds. Loading is mass per entire
horizontal ground-cell area. The consumed fraction describes fuel consumed
in the modeled phase, and effective released heat is a supplied value; neither
is automatically corrected for moisture. Define the conditional released
energy density and conditional heat-release-rate density by
\begin{align}
 e_j&=w_jf_j H_{c,j}
       &&[\mathrm{J\,m^{-2}}],\\
 Q_j&=\frac{e_j}{T_j}=\frac{w_jf_j H_{c,j}}{T_j}
       &&[\mathrm{W\,m^{-2}}].
 \label{eq:heat-hrr}
\end{align}
For an active cell, the modeled total heat-release rate is
$Q_jA_{\rm cell}$ watts. With prescribed radiative fraction
$\chi_j\in[0,1]$, the active-cell radiant rate is
$\chi_j Q_j A_{\rm cell}$ watts. The radiative fraction partitions released
power; it is not receiver absorptivity or atmospheric transmissivity.
The \texttt{CombustionParameters} object evaluates
\eqref{eq:heat-hrr}. Its fields can broadcast during construction, but the
radiation model ultimately requires $Q_j$ and $\chi_j$ to be either scalars
or spatial maps with exactly the probability map's $[H,W]$ shape.

\subsection{From a first-burning probability to active occupancy}
\label{sec:heat-timing}

The maintained checkpoint predicts first new cumulative burning during an
exact one-hour interval, whose duration is
$D=3600\ \mathrm{s}$. Let $N_j$ be that event, and define its
\emph{unconditional with respect to initial activity status} forecast
probability, conditional on the supplied information $\mathcal F_0$, by
\begin{equation}
 q_j=\mathbb P(N_j\mid\mathcal F_0).
\end{equation}
Let $U_j$ be the arrival time, relative to the cutoff, conditional on $N_j$.
The temporal bridge adds the assumptions
\begin{equation}
 U_j\mid(N_j,\mathcal F_0)\sim\operatorname{Uniform}(0,D),
 \qquad
 Z_j^{\rm new}(t)=\Ind_{N_j}\Ind_{\{U_j\le t<U_j+T_j\}}.
 \label{eq:heat-uniform-arrival}
\end{equation}
The conditional law is required only when $q_j>0$. When $q_j=0$,
new activity is identically zero. Extend $U_j$ arbitrarily off $N_j$, where
its value cannot affect the indicator. A product of $q_j$ and a conditional
expectation given $N_j$ is interpreted as zero in the $q_j=0$ branch.
The assigned model or scenario probability $q_j$ is not an independently
established true physical arrival probability. All expectation identities
below are exact under the assigned probability law and supplied parameters.
The trained predictor supplies $q_j$; it does not supply the uniform
conditional arrival law or $T_j$. No assumption about dependence between
arrival times in different cells is required for the expectations below.

For $0\le t\le D$, the conditional interval of arrivals that remain active
at $t$ is $(\max(0,t-T_j),t]$. Its length is $\min(t,T_j)$, except for
endpoints of zero probability under the continuous uniform law. Consequently,
\begin{equation}
 p_j^{\rm new}(t)
 =\mathbb P(Z_j^{\rm new}(t)=1\mid\mathcal F_0)
 =q_j\frac{\min(t,T_j)}{D}.
 \label{eq:heat-instantaneous-occupancy}
\end{equation}
Define the time-mean active probability and endpoint active probability by
\begin{equation}
 \bar p_j^{\rm new}=\frac1D\int_0^D p_j^{\rm new}(t)\,dt,
 \qquad p_{j,D}^{\rm new}=p_j^{\rm new}(D).
\end{equation}
If $T_j\le D$, split the integral at $T_j$:
\begin{align}
 \bar p_j^{\rm new}
 &=\frac{q_j}{D^2}
      \left(\int_0^{T_j}t\,dt+\int_{T_j}^D T_j\,dt\right)\\
 &=q_j\left(\frac{T_j}{D}-\frac{T_j^2}{2D^2}\right).
\end{align}
If $T_j\ge D$, the integrand is $q_jt/D$ over the whole forecast interval,
so $\bar p_j^{\rm new}=q_j/2$. Writing
$\rho_j=\min(T_j,D)/D\in(0,1]$ gives both cases in one expression:
\begin{equation}
 \boxed{\bar p_j^{\rm new}=q_j(\rho_j-\rho_j^2/2),
 \qquad p_{j,D}^{\rm new}=q_j\rho_j.}
 \label{eq:heat-occupancy-bridge}
\end{equation}
Equivalently, a cell arriving at time $u$ burns for
$\min(T_j,D-u)$ seconds within $[0,D]$, and
$\bar p_j^{\rm new}=q_j\E[\min(T_j,D-U_j)\mid N_j,\mathcal F_0]/D$.
The corresponding expected radiant energy emitted by a newly igniting cell
\emph{within the forecast interval}, before atmospheric loss, is
\begin{equation}
 A_{\rm cell}\chi_j Q_jD\bar p_j^{\rm new}
 =A_{\rm cell}\chi_j e_jq_j
 \begin{cases}
  1-T_j/(2D),&T_j\le D,\\
  D/(2T_j),&T_j\ge D.
 \end{cases}
\end{equation}
This generally differs from $A_{\rm cell}\chi_j e_jq_j$, because a cell
arriving late can continue burning beyond the forecast interval. Thus the
time-mean occupancy is the expected fraction of the forecast interval spent
active; it is not an instantaneous state at the midpoint.
For example, $T_j=300\ \mathrm{s}$ and $D=3600\ \mathrm{s}$ give
$\rho_j=1/12$, $\bar p_j^{\rm new}=(23/288)q_j$, and
$p_{j,D}^{\rm new}=q_j/12$.

The probability $q_j$ alone does not determine active occupancy. Arrivals
concentrated near zero and arrivals concentrated near $D$ can have the same
$q_j$ and different interval dose and endpoint activity. Equation
\eqref{eq:heat-occupancy-bridge} is valid for the implemented additional
arrival scenario. The generic temporal utilities allow any positive numeric
$D$, while \texttt{WildfireHeatModel} enforces $D=3600$ because its maintained
predictors were fitted on exact one-hour pairs.

\subsection{Explicit contribution from fire already active at the cutoff}
\label{sec:heat-present}

Cumulative burning by time zero does not establish current combustion.
Existing fire is included only when the caller supplies both its
initial active probability $p_{j,0}^{\rm old}\in[0,1]$ and its deterministic
remaining active duration $R_j\ge0$ seconds conditional on that initial
state. The modeled old-fire contribution has active interval $[0,R_j)$:
\begin{equation}
 p_j^{\rm old}(t)=p_{j,0}^{\rm old}\Ind_{\{0\le t<R_j\}},
 \qquad
 \bar p_j^{\rm old}=p_{j,0}^{\rm old}\frac{\min(R_j,D)}{D},
 \qquad
 p_{j,D}^{\rm old}=p_{j,0}^{\rm old}\Ind_{\{R_j>D\}}.
 \label{eq:heat-initial-occupancy}
\end{equation}
The strict endpoint inequality matters: $R_j=D$ gives extinction exactly
at the endpoint and therefore zero endpoint activity. The API permits
$R_j=0$, which contributes zero throughout the modeled interval. The same
prescribed conditional source density $\chi_jQ_j$ is applied to present
fire as to new fire. The numerical checks require $R_j\ge0$ but do not
require $R_j\le T_j$ or verify the remaining fuel inventory; consistency
of an old-fire phase and its remaining duration is a caller-supplied
physical assumption.

The pipeline requires positive initial active probability only in explicitly
known already-burned cells. It sets the first-new-burning probability to zero
in those cells and marks that zero valid even if the learned neighborhood is
unsupported. Initial-active sources and possible first-new-burning sources
therefore have disjoint cell support in this implementation. Then
\begin{equation}
 \bar p_j=\bar p_j^{\rm new}+\bar p_j^{\rm old},
 \qquad p_{j,D}=p_{j,D}^{\rm new}+p_{j,D}^{\rm old}.
\end{equation}
These additions use unconditional event probabilities. More general
same-cell combinations would require disjoint events and consistent
conditioning; a probability conditional on initially being unburned would
first need multiplication by the probability of that conditioning event.
The implementation rejects overlapping positive initial activity and
possible new burning instead of silently treating them as independent.
If initial inputs are absent, $p_j^{\rm old}(t)$ is set to zero as an explicit
\emph{new-ignition-only scope}; no inference of old-fire age is made.

\subsection{Radiation geometry and the exact finite-cell transfer}
\label{sec:heat-geometry}

All source cells are represented by horizontal emitting sheets at common
height $z_s$, and all receivers lie at common height $z_r$ in the same datum.
The required separation is $\ell=z_s-z_r>0$. The receivers are upward-facing
horizontal surfaces with unit normal $\mathbf n=(0,0,1)$. The common datum
can be shifted, even to negative heights; only $\ell$ affects this model.
Terrain elevation, slope, and wind may affect the learned probability model,
but they do not tilt a receiver, bend a flame, or change the sheet heights in
the radiation operator.

For a horizontal receiver point $\mathbf x_i$ and horizontal source location
$\mathbf y\in S_j$, where $S_j$ is cell $j$'s footprint, define
\begin{equation}
 r(\mathbf y,\mathbf x_i)=
    \sqrt{\|\mathbf y-\mathbf x_i\|^2+\ell^2},
 \qquad
 \cos\theta=\frac{\ell}{r(\mathbf y,\mathbf x_i)}.
\end{equation}
The angle $\theta$ is between the upward receiver normal and the direction
from the receiver toward the source. An infinitesimal horizontal source
area $dA$ has isotropic radiant power density $\chi_jQ_j$, giving conditional
radiant rate $\chi_jQ_j\,dA$. Isotropic emission spreads that rate over
$4\pi r^2$, and projecting it onto the receiver contributes the factor
$\cos\theta$. For fixed scalar atmospheric transmissivity
$\mathcal T\in[0,1]$, the resulting transfer is
\begin{equation}
 K_{ij}=\frac1{4\pi}\int_{S_j}
        \frac{\ell}{(\|\mathbf y-\mathbf x_i\|^2+\ell^2)^{3/2}}\,dA
       =\frac{\Omega_{ij}}{4\pi}.
 \label{eq:heat-transfer}
\end{equation}
Here $\Omega_{ij}$ is the solid angle in steradians subtended by the
horizontal parallelogram at the receiver. The integral is dimensionless:
$\ell/r^3$ has unit $\mathrm{m^{-2}}$ and $dA$ has unit $\mathrm{m^2}$.
The equality to solid angle uses the horizontal source sheet and parallel
horizontal receiver normal. There is no additional Lambertian source-normal
factor: the infinitesimal emitters are isotropic, rather than a solid-flame
surface with a directional emission law. Arbitrary receiver normals or
vertical receiving walls are not supported by the current API.\cite{heatrefs}

For completeness, the implemented solid-angle calculation is as follows.
Write the four source-corner vectors from the receiver as
\begin{equation}
 \mathbf v_k=(\mathbf x_j-\mathbf x_i
             +\tfrac12(\alpha_k\mathbf u+\beta_k\mathbf v),\;\ell),
 \quad
 (\alpha_k,\beta_k)\in\{(-1,-1),(1,-1),(1,1),(-1,1)\},
\end{equation}
in the displayed order $k=0,1,2,3$. For a triangle with corner indices
$(a_1,a_2,a_3)$, put $n_k=\|\mathbf v_k\|$ and define
\begin{align}
 D_{a_1a_2a_3}
 &=n_{a_1}n_{a_2}n_{a_3}
 +(\mathbf v_{a_1}\cdot\mathbf v_{a_2})n_{a_3}
 +(\mathbf v_{a_2}\cdot\mathbf v_{a_3})n_{a_1}
 +(\mathbf v_{a_3}\cdot\mathbf v_{a_1})n_{a_2},\\
 \Omega_{a_1a_2a_3}
 &=2\operatorname{atan2}(\ell A_{\rm cell},D_{a_1a_2a_3}).
\end{align}
The magnitude of the triple product for either triangular half of the cell
is $\ell A_{\rm cell}$. Splitting along diagonal $(0,2)$ gives
\begin{equation}
 \Omega_{ij}=\Omega_{012}+\Omega_{023},
 \qquad K_{ij}=\frac{\Omega_{012}+\Omega_{023}}{4\pi}.
\end{equation}
The routine constrains roundoff to $0\le K_{ij}\le1/2$. The use of finite
cells and strictly positive separation replaces the singular same-location
point-source approximation with a finite integral. There is no tunable
minimum distance, added squared-distance constant, distance cutoff, or
normalized spatial kernel. At a receiver under the interior of a cell,
$\ell\downarrow0$ gives $\Omega_{ij}\to2\pi$ and $K_{ij}\to1/2$.

In the far field, when a cell is small compared with its distance $r_{ij}$
from the receiver, the transfer approaches
\begin{equation}
 K_{ij}\simeq\frac{A_{\rm cell}\ell}{4\pi r_{ij}^3},
 \qquad
 F_{ij}\simeq\frac{\mathcal T P_j^{\rm rad}}{4\pi r_{ij}^2}\frac{\ell}{r_{ij}},
\end{equation}
where $P_j^{\rm rad}=\chi_j Q_jA_{\rm cell}$ is active-cell radiant power. The factor
$\ell/r_{ij}$ is required by the upward receiver orientation. As a separate
consistency check, the infinite horizontal receiving plane intercepts
$\mathcal T/2$ of total isotropic source power; a uniform infinite source sheet
of density $s$ gives flux $\mathcal T s/2$. Neither identity permits identifying
the sum of finite-grid receiver samples with emitted source power.

Atmospheric transmissivity is a single supplied path-independent scalar.
The implementation multiplies the final irradiance by $\mathcal T$, without
changing the emitted radiant source power. It does not compute path length,
humidity or smoke absorption, occlusion, shielding, or source temperature.

\subsection{Why marginal probabilities suffice for the expected flux}
\label{sec:heat-expectation}

Condition on $\mathcal F_0$ and on the supplied fixed physical parameters and
geometry. The instantaneous incident radiant flux is the modeled random sum
\begin{equation}
 F_i(t)=\mathcal T\sum_{j\in\mathcal J}K_{ij}\chi_jQ_jZ_j(t)
 \qquad[\mathrm{W\,m^{-2}}].
\end{equation}
Because the source domain is finite, expectation can be taken term by term:
\begin{align}
 \E[F_i(t)\mid\mathcal F_0]
 &=\mathcal T\sum_{j\in\mathcal J}K_{ij}\chi_jQ_j
       \E[Z_j(t)\mid\mathcal F_0]\\
 &=\boxed{\mathcal T\sum_{j\in\mathcal J}K_{ij}\chi_jQ_jp_j(t)}.
 \label{eq:heat-linear-expectation}
\end{align}
This calculation requires no independence between the $Z_j(t)$. For
example, perfectly correlated burning in two cells and mutually exclusive
burning in those cells can yield identical marginal probabilities and thus
identical expected fluxes, even though their fluctuation and exceedance
probabilities differ. The present API computes an expectation, not those
joint-distribution-dependent quantities.

The source strengths are prescribed conditional strengths, not random
unconditional averages. If actual active source density $\Xi_j$ in $\mathrm{W\,m^{-2}}$ were random,
then in general
\begin{equation}
 \E[Z_j \Xi_j\mid\mathcal F_0]
 =p_j\E[\Xi_j\mid Z_j=1,\mathcal F_0],
\end{equation}
whenever the conditional expectation is defined. It need not equal
$p_j\E[\Xi_j\mid\mathcal F_0]$. Treating $\chi_jQ_j$ as fixed avoids
silently imposing that incorrect factorization. Uncertainty in fuel,
consumption, duration, radiative fraction, geometry, and transmission is not
marginalized by this package.

The expected radiant source power returned by the radiation-only model is
\begin{equation}
 \E[P_{\rm radiant}(t)\mid\mathcal F_0]
 =A_{\rm cell}\sum_{j\in\mathcal J}\chi_jQ_jp_j(t)
 \qquad[\mathrm W].
 \label{eq:heat-source-power}
\end{equation}
This emitted source power excludes atmospheric loss. It is distinct from
incident receiver flux or power.

Since the operator and source strengths are time independent, integration
in time commutes with the finite sum. The interval-mean expected flux,
expected dose per receiver area, and endpoint expected flux are respectively
\begin{align}
 \bar F_i&=\mathcal T\sum_j K_{ij}\chi_jQ_j\bar p_j
           &&[\mathrm{W\,m^{-2}}],\\
 \mathcal E_i&=\int_0^D\E[F_i(t)\mid\mathcal F_0]\,dt
             =D\bar F_i
           &&[\mathrm{J\,m^{-2}}],\\
 F_{i,D}:=\E[F_i(D)\mid\mathcal F_0]&=\mathcal T\sum_j K_{ij}\chi_jQ_jp_{j,D}
           &&[\mathrm{W\,m^{-2}}].
 \label{eq:heat-mean-dose}
\end{align}
Thus multiplying the exact interval mean by $D$ is valid even when
instantaneous occupancy changes during the hour. Multiplying endpoint flux
by $D$ would require a different constant-flux assumption and is not this
pipeline's dose calculation.

\subsection{Evaluation at raster centres or arbitrary points}
\label{sec:heat-evaluation}

For an aligned affine source and receiver grid with common heights,
$K_{ij}$ depends only on row and column displacement. The source matrix is
$s_j=p_j\chi_jQ_j$. The implementation builds the transfer array with
shape $[2H-1,2W-1]$ and evaluates the linear convolution of that array with
$s$. It pads both arrays to full convolution shape
$[3H-2,3W-2]$ before transforming them, multiplies their Fourier transforms,
inverts the transform, and extracts rows $H-1$ through $2H-2$ and columns
$W-1$ through $2W-2$. The Fourier transform accelerates this sum; it does
not change the radiation law or approximate the finite-cell solid angles.
Padding prevents a source near one edge from wrapping periodically to the
opposite edge. Source values outside the supplied raster are excluded by a
finite-domain assumption; this padding does not establish that unavailable
source data are physically zero. Small negative floating-point values after
the inverse transform are replaced by zero because the exact mathematical
sum is nonnegative. Input probabilities are validated rather than clipped.

Alternatively, \texttt{predict\_at\_points} accepts an $[N,2]$ array of
arbitrary horizontal projected receiver coordinates and sums the same exact
cell transfers directly. Its default target and source chunk sizes are
128 and 1024, respectively. Chunking bounds temporary memory usage and has
no physical meaning. Zero-strength cells can be skipped in this direct sum.
Every receiver still lies on the configured common horizontal plane.

The raster output contains flux values at receiver cell \emph{centres}, not
averages of flux over receiver cells. If a small receiving area
$A_i^{\rm recv}\ge0$ is assumed to experience uniform flux, then its
incident radiant power is $\bar F_iA_i^{\rm recv}$ watts, and its
interval incident energy is $D\bar F_iA_i^{\rm recv}$ joules. A large
receiving patch $S$ requires the area integral
$\int_S\bar F(\mathbf x)\,dA$, or a sufficiently resolved numerical
approximation. Multiplication of one point sample by the whole source-cell
area is not generally that integral. None of these quantities is absorbed
heat; that would additionally require an appropriate receiver absorption
model.

For a separate sampled flux time series at strictly increasing seconds
$t_0<\cdots<t_M$, \path{integrate_heat_flux} uses the trapezoidal rule
\begin{equation}
 \mathcal E_i^{\rm trap}
 =\sum_{m=0}^{M-1}\frac{F_i(t_m)+F_i(t_{m+1})}{2}(t_{m+1}-t_m).
\end{equation}
This assumes linear interpolation between supplied samples. It does not
reconstruct missing intervals or rapid unsampled variations. For
piecewise-constant interval forecasts, the corresponding calculation is
$\sum_m\bar F_{i,m}\Delta t_m$.

\subsection{Unknown sources, explicit estimates, and monotone bounds}
\label{sec:heat-bounds}

The radiation-only model rejects missing, masked, nonfinite, negative, or
above-one active probabilities. The integrated pipeline can preserve
unsupported learned probabilities using an explicit Boolean validity mask.
Let $V\subseteq\mathcal J$ be the set of resolved $q_j$, after setting
known already-burned first-event probabilities to zero. A valid entry must
be finite and in $[0,1]$. An invalid entry is retained as NaN even when its
raw stored number is finite; validity is part of the data contract.

An optional externally justified probability map can be supplied for
unsupported entries. It changes only those entries and is recorded using
an externally-supplied-probability mask. The pipeline also retains the
original validity mask, so a supplied estimate cannot be mistaken for the
checkpoint's supported prediction.

For entries still unresolved, let
\begin{equation}
 q_j^-=
 \begin{cases}q_j,&j\in V,\\0,&j\notin V,\end{cases}
 \qquad
 q_j^+=
 \begin{cases}q_j,&j\in V,\\1,&j\notin V.\end{cases}
\end{equation}
Put $g_j=\rho_j-\rho_j^2/2$ and $b_j=\bar p_j^{\rm old}$.
For any complete map $q^*$ consistent with the known entries and
$0\le q_j^*\le1$, define
\begin{equation}
 \bar F_i(q^*)=\mathcal T\sum_jK_{ij}\chi_jQ_j(g_jq_j^*+b_j).
\end{equation}
Every coefficient in this sum is nonnegative. Therefore the implemented
interval-mean bounds satisfy
\begin{equation}
 \boxed{\bar F_i^-:=\bar F_i(q^-)
 \ \le\ \bar F_i(q^*)\ \le\ \bar F_i(q^+)=:\bar F_i^+.}
 \label{eq:heat-bounds}
\end{equation}
Endpoint bounds replace $g_j$ by $\rho_j$ and $b_j$ by
$p_{j,D}^{\rm old}$. Dose bounds are $D\bar F_i^-$ and $D\bar F_i^+$.
The lower and upper substitutions are mathematical uncertainty endpoints,
not observations or default assumptions about missing fire.
The finite known-source contribution is $\bar F_i^-$, including any
explicitly supplied present-active sources. It is not the unknown full
expectation.

With $\ell>0$ every nonempty source cell has positive geometric transfer to
every finite receiver point. Consequently, an unresolved $q_j$ for which
$\mathcal T\chi_jQ_j>0$ can affect the entire receiver field, not merely its
corresponding pixel. The pipeline marks the full expected mean flux and
expected dose as NaN throughout the receiver raster if any such emitting
source remains unresolved, and sets \texttt{complete=False}. The bounds
and known-source contribution remain finite. If every unresolved cell has
zero source density, zero radiative fraction, or $\mathcal T=0$, the full flux is
nevertheless determined, so \texttt{complete=True} is possible even with
unresolved probability entries. The occupancy entries themselves remain
unknown.

These are bounds for fixed supplied combustion, timing, transmission,
geometry, and finite source domain. They do not bound unmodeled transport,
outside-domain fires, parameter error, active flames omitted from the
present-fire scope, or errors in the learned probability model.

\subsection{Surface-flame scenarios and position selection}
\label{sec:heat-flames}

The point-residence construction and its reference-fuel assumptions are
described in Section~\ref{sec:residence}. A residence duration describes
flaming at a representative fuel location, not time to cross the entire
375 m raster cell or time until all smoldering and trailing combustion ends.

For a supplied known or estimated arrival $a_j$ seconds and query time
$t$ seconds sharing a common origin, \texttt{active\_flames\_from\_arrival}
defines the scenario indicator and remaining duration by
\begin{align}
 z_j^{\rm scenario}(t)&=\Ind_{\{a_j\le t<a_j+\tau_j^{\mathrm{fl}}\}},\\
 R_j^{\rm scenario}(t)&=
 \begin{cases}\tau_j^{\mathrm{fl}}-(t-a_j),&a_j\le t<a_j+\tau_j^{\mathrm{fl}},\\0,&\text{otherwise.}\end{cases}
\end{align}
If a reconstructed arrival $\widehat a_j$ is recorded in hours from the
event origin and the query uses that same origin, its conversion is
$a_j=3600\widehat a_j$ seconds. If time zero instead denotes the forecast
cutoff, subtract the cutoff elapsed time from that converted arrival before
forming $t-a_j$. Mixing event-relative hours and cutoff-relative seconds
would give an incorrect activity indicator.
The implementation compares elapsed age $t-a_j$ with $\tau_j^{\mathrm{fl}}$, avoiding an
unnecessary overflow from forming $a_j+\tau_j^{\mathrm{fl}}$. It permits negative arrival
and query times when their common origin warrants them. At exactly
$a_j+\tau_j^{\mathrm{fl}}$, the location is inactive. Invalid arrivals remain NaN and
unclassified. An estimated reconstructed arrival yields a nominal scenario;
this function does not integrate arrival uncertainty or identify measured
flame positions.

Define $\rho_j^{\mathrm{fl}}=\min(\tau_j^{\mathrm{fl}},D)/D$. The
flame-specific mean is $q_j[\rho_j^{\mathrm{fl}}-\tfrac12(\rho_j^{\mathrm{fl}})^2]$,
and the endpoint is $q_j\rho_j^{\mathrm{fl}}$. These specialize the timing
bridge to surface-flame residence; they agree with the heat-emission bridge
only when the caller explicitly sets $T_j=\tau_j^{\mathrm{fl}}$.
For uncertain new arrivals, \texttt{forecast\_active\_flames} returns
$q_j$, the mean flame occupancy specialized above, and an
endpoint map containing $q_j\rho_j^{\mathrm{fl}}$. Existing flames are excluded from
this forecast unless supplied separately through the integrated heat
pipeline. Unknown probabilities remain NaN. A requested selection cutoff
$\alpha\in(0,1]$ returns exactly the valid cells whose instantaneous
probability is at least $\alpha$, together with their zero-based row and
column indices and projected cell-centre coordinates. There is no implicit
selection threshold. Thresholding is a display or candidate-selection rule;
it does not establish detection accuracy or calibration against active-flame
observations. Empty selections remain usable arrays of shape $[0,2]$ for
indices and coordinates and $[0]$ for probabilities.

\subsection{The composed checkpoint-to-heat procedure}
\label{sec:heat-pipeline}

The public \texttt{WildfireHeatModel} executes the following sequence for
one spatial snapshot.
\begin{enumerate}
 \item Load the saved checkpoint, identify either the maintained direct
 environment-conditioned convolution or survival-intensity convolution
 architecture, and reuse its saved channel schema, categorical vocabulary,
 fitted weights, and training-only normalization statistics. The integration
 does not train or fit preprocessing. Checkpoint path, SHA-256, architecture,
 selected epoch, channel information, and device are recorded. CPU is the
 default device.
 \item Establish a common projected affine grid, complete combustion scenario,
 radiative fraction, source/receiver heights, and transmission scalar. If a
 \texttt{FlameResidenceTime} object is also supplied, its duration must equal
 \texttt{combustion.burning\_duration\_s} exactly at every represented cell.
 The same duration then governs active occupancy and the denominator in
 conditional heat release. Consumed fraction must concern that modeled
 flaming phase.
 \item Pass the current cumulative burned state, current environmental tensor,
 exact channel order, and separate validity masks to the saved predictor.
 Use a one-hour interval, \path{dt_hours=1.0}, and an unknown boundary,
 \path{boundary="unknown"}. Only a
 single $[H,W]$ snapshot, possibly with singleton leading dimensions, is
 accepted. For sample-based heat prediction, the affine coefficients must
 agree with the configured radiation grid to absolute tolerance $10^{-9}$
 and zero relative tolerance.
 \item Take the predictor's evaluated probability and validity outputs. For a
 direct convolution this means its bounded evaluation probability, not its
 unbounded signed raw score. Resolve first-new-burning probability to zero
 in known already-burned cells. Retain unsupported entries and separately
 record any caller-supplied estimates.
 \item Convert new-event probabilities to mean and endpoint active occupancy
 using the uniform-arrival and constant-duration scenario. Add only explicitly
 supplied initial-active contributions with known already-burned support
 and remaining durations.
 \item Apply the fixed radiation operator to lower and upper mean occupancies
 and lower and upper endpoint occupancies. Return a complete expected mean
 only when unknown emitting contributions are absent. Otherwise return
 the finite known-source contribution and monotone bounds, preserving NaN
 for the unresolved full expectation.
 \item Multiply mean flux and its bounds by $3600\ \mathrm{s}$ to obtain
 interval radiant dose per receiver area. Preserve probability masks,
 externally supplied masks, temporal and physical assumptions, coverage
 counts, completeness, and provenance in the result.
\end{enumerate}
The sample adapter does not consume arrival surfaces, arrival uncertainty,
or next-hour labels to run the checkpoint. Nevertheless, a current burned
state reconstructed retrospectively using an arrival surface and final
extent remains a retrospective input. Restricting environmental intervals
to end before the cutoff does not convert that reconstruction into an
operationally observed current fire state.

\subsection{Exact scope of the saved demonstrations and verification}
\label{sec:heat-demo}

The saved \path{integrated_example.py} demonstration uses the Fish Creek
event (identifier \path{WY-BTF-002416_FishCreek}). Its cutoff is
\path{2024-08-23T00:00:00Z}; it uses the signed direct $5\times5$ checkpoint,
and a $103\times73$ grid in EPSG:5070 with 375 m cells. Its specified
parameters are $H_{c,j}=18{,}000{,}000\ \mathrm{J\,kg^{-1}}$,
$T_j=300\ \mathrm{s}$, $f_j=0.10$, $\chi_j=0.30$,
$z_s=5\ \mathrm m$, $z_r=0\ \mathrm m$, and $\mathcal T=1$.
The ecological loading adapter uses an explicit
$10\ \mathrm{kg\,m^{-2}}$ loading scenario for unmapped fractions.
Those are illustrative scenario values, not fitted heat-flux parameters.
The saved reference-versus-uniform comparison keeps the same probability
map and compares ecological loading with constant
$10\ \mathrm{kg\,m^{-2}}$ loading. It establishes sensitivity to that
input, not an accuracy improvement. The historical summary records
225 resolved probabilities, 7294 unresolved probabilities, and
\path{complete_expected_flux_available=False}; the finite plotted
known-source contribution must therefore not be relabeled as full expected
heat.\cite{codeheat}

The separate \texttt{active\_flame\_example.py} uses the same default event,
cutoff, and checkpoint but constructs a spatial reference-fuel residence
map, rather than using the heat demonstration's constant 300 s duration.
Its default unsupported-code SAV fallback is explicitly
$2000\ \mathrm{ft^{-1}}$, and the default candidate cutoff is $0.001$.
The saved example preserves a zero nominal-current-flame selection and
zero forecast candidates at that cutoff, while also exporting all positive
endpoint probabilities. A zero selected count is retained as a result of
the specified scenario and threshold. It is not evidence that no flames
physically existed. The reconstructed-arrival uncertainty is much larger
than the point-scale residence duration, and is not integrated into the
nominal indicator.

The current source tests in \path{tests/heat/test_model.py},
\path{tests/heat/test_temporal.py}, \path{tests/heat/test_pipeline.py},
\path{tests/heat/test_flames.py}, and
\path{tests/heat/test_flame_integration.py} specify independent mathematical
checks: source-area Gauss--Legendre integration, analytic rectangle solid
angles, the orientation-corrected far-field limit, half-power interception
by an infinite receiver plane, equal expectations for distinct dependence
structures with identical marginals, affine rotation/shear, lack of periodic
wrap, temporal indicator quadrature, endpoint extinction conventions,
unknown-source bounds, units, and preservation of saved preprocessing and
masks. These checks verify the stated mathematical model and software
contracts. Historical validation-report files record earlier runs, sometimes
using paths predating test relocation. They are not substitutes for a fresh
current test result and do not constitute measured-wildfire heat or
active-flame calibration.

\section{Step-by-step worked calculations and public interfaces}
\label{sec:worked}
\subsection{A complete one-cell calculation with an analytic answer}
This example isolates the timing and radiation equations. It is synthetic and
is not a fitted fire scenario. Consider one square $10\ \mathrm m$ cell,
with its receiver at the horizontal center, source height 5 m, receiver height
0 m, and transmissivity $\mathcal T=1$. Supply $q=0.8$ for first new burning
in the next $D=3600$ seconds, no initial active fire, and
\begin{equation}
w=1\ \mathrm{kg\,m^{-2}},\quad f=1,\quad
H_c=18{,}000{,}000\ \mathrm{J\,kg^{-1}},\quad
T=120\ \mathrm s,\quad\chi=0.3.
\end{equation}
Each step is now fully determined:
\begin{enumerate}
\item The cell area is $100\ \mathrm{m^2}$, but dry fuel loading is already
mass per unit whole-cell ground area. Released energy density is
$e=w f H_c=18{,}000{,}000\ \mathrm{J\,m^{-2}}$.
\item Conditional heat-release density is $Q=e/T=150{,}000\ \mathrm{W\,m^{-2}}$.
An active cell therefore releases $15{,}000{,}000$ W in total, of which
$4{,}500{,}000$ W is radiative under $\chi=0.3$.
\item The capped duration ratio is $r=T/D=1/30$. Under the uniform-arrival
law, the mean and endpoint activity probabilities are
\begin{equation}
\overline p=0.8\left(\frac1{30}-\frac1{1800}\right)
=\frac{59}{2250}\simeq0.02622222,
\qquad p(D)=\frac{0.8}{30}=\frac2{75}\simeq0.02666667.
\end{equation}
\item For a centered rectangle with half-width $a=5$ m, half-height $b=5$ m,
and vertical separation $\ell=5$ m, its solid angle is
\begin{equation}
\Omega=4\arctan\left(\frac{ab}{\ell\sqrt{\ell^2+a^2+b^2}}\right)
=4\arctan(1/\sqrt3)=\frac{2\pi}{3}.
\end{equation}
Thus the finite-cell transfer $\Omega/(4\pi)$ is $1/6$.
\item Mean incident radiant flux is
\begin{equation}
\overline F=\mathcal T\,\overline p\,\chi Q\frac16
=\frac{590}{3}\ \mathrm{W\,m^{-2}}
\simeq196.6666667\ \mathrm{W\,m^{-2}}.
\end{equation}
Endpoint flux is $F_D=p(D)\chi Q/6=200\ \mathrm{W\,m^{-2}}$.
Neither equals the next-hour transition probability multiplied directly by
the conditional active-cell flux.
\item Expected interval incident energy is
\begin{equation}
\mathcal E=D\overline F=708{,}000\ \mathrm{J\,m^{-2}}.
\end{equation}
This is a point receiver's dose per receiver area, not total energy received
by the whole ground cell.
\item All source probabilities and physical parameters are supplied and finite,
so this single-source calculation is complete. A missing $q$ would instead
produce a range by substituting 0 and 1, under the same fixed scenario.
\end{enumerate}

The public objects evaluate this example without any wildfire checkpoint:
\begin{lstlisting}[language=Python]
import numpy as np
from expected_heat_flux import (
    AffineGrid, CombustionParameters, ExpectedHeatFluxModel,
    ignition_to_active_probability,
)

combustion = CombustionParameters(
    fuel_load_kg_m2=1.0,
    heat_of_combustion_j_kg=18_000_000.0,
    burning_duration_s=120.0,
    consumed_fraction=1.0,
)
timing = ignition_to_active_probability(
    np.array([[0.8]]), 120.0, 3600.0,
)
radiation = ExpectedHeatFluxModel(
    grid=AffineGrid.from_cell_size(10.0),
    heat_release_rate_density_w_m2=(
        combustion.heat_release_rate_density_w_m2
    ),
    radiative_fraction=0.3,
    emission_height_m=5.0,
    receiver_height_m=0.0,
    atmospheric_transmissivity=1.0,
)
mean = radiation.predict(timing.mean_active_probability)
endpoint = radiation.predict(timing.endpoint_active_probability)
mean_flux = mean.incident_heat_flux_w_m2[0, 0]
endpoint_flux = endpoint.incident_heat_flux_w_m2[0, 0]
np.testing.assert_allclose(mean_flux, 590.0 / 3.0)
np.testing.assert_allclose(endpoint_flux, 200.0)
np.testing.assert_allclose(mean_flux * 3600.0, 708_000.0)
\end{lstlisting}
The code names, call signatures, and assertions above were executed against
the current package during document preparation. Only the array's active probability is
passed to the independent radiation object; no event-transition probability
is implicitly reinterpreted inside that object.

\subsection{A residence-based endpoint example}
With characteristic SAV $2000\ \mathrm{ft^{-1}}$, the correctly converted SI
value is $2000/0.3048\ \mathrm{m^{-1}}$. Anderson--Andrews residence is
$\tau^{\mathrm{fl}}=11.52$ s. For $q=0.8$ and the same one-hour horizon,
\begin{equation}
p^{\mathrm{fl}}(D)=0.8\frac{11.52}{3600}=0.00256,
\qquad
\overline p^{\mathrm{fl}}
=0.8\left(0.0032-\frac{0.0032^2}{2}\right)=0.002555904.
\end{equation}
The calculation does not reduce the probability of burning sometime during
the hour; that remains 0.8. It changes the event whose probability is being
asked for. Endpoint flame activity requires ignition sufficiently close to
the endpoint to remain active under the supplied short residence duration.

\subsection{Loading exactly the saved model's input contract}
Run the following from the directory containing the root \path{README.md}
and \path{pyproject.toml}, using the project's Python 3.12 environment. Its
dependency pins include NumPy 1.26.4 and PyTorch 2.5.1. The source/data paths
below are repository-relative for portability. An installed code-only wheel
requires explicit access to the separately retained data and checkpoint.
\begin{lstlisting}[language=Python]
from pathlib import Path
from wildfire_burned_area_hourly import load_model

root = Path.cwd()
checkpoint = (
    root / "wildfire_burned_area_hourly/probability_model"
    / "training_runs/20261007_1h_direct_convolution_signed"
    / "5x5/fitted/checkpoint.pt"
)
model = load_model(checkpoint, device="cpu")
with model.load_data(
    root / "wildfire_burned_area_hourly/environmental_data",
    events=["WY-BTF-002416_FishCreek"],
    sequence_length=1,
    time_range=("2024-08-23T00:00:00Z", "2024-08-23T00:00:00Z"),
) as dataset:
    assert len(dataset) == 1
    sample = dataset[0]

prediction = model.predict_sample(sample)
assert tuple(sample["channel_names"]) == model.channel_names
assert tuple(prediction.probability.shape) == (1, 1, 103, 73)
print(int(prediction.valid.sum()))  # 225 in the verified example
\end{lstlisting}
The example's 7519 raster locations comprise 225 valid probabilities and 7294
unknown probabilities. The CPU run during authorship produced maximum known
probability $0.09081834554672241$. These counts describe prediction coverage,
not accuracy; the example itself does not consume future labels for inference.
The checkpoint predicts raw and clipped values separately. Use
\path{prediction.probability} and \path{prediction.valid} for probability
calculations, and retain \path{prediction.raw_score} for raw-score analysis.

To connect this snapshot to heat, construct a metric \path{AffineGrid} from
the sample's six affine coefficients, supply complete combustion inputs and
geometry, and call \path{model.heat_model(...).predict_sample(sample)}.
The result's \path{complete} field determines whether the full supplied-domain
expected heat is known. Export arrays, their masks, affine transform, CRS,
timestamp, physical assumptions, and provenance together. A local valid
probability does not suffice to establish a locally complete heat prediction
when radiation from another source is unknown.

\section{What the calculation establishes, and what remains unspecified}
\label{sec:limits}
\subsection{Learned quantities, supplied quantities, and absent quantities}
\begin{longtable}{@{}P{0.25\linewidth}P{0.34\linewidth}P{0.33\linewidth}@{}}
\toprule Stage & Computed or learned here & Additional information not supplied by that stage\\\midrule
\endfirsthead
\toprule Stage & Computed or learned here & Additional information not supplied by that stage\\\midrule
\endhead
Arrival product & Retrospective location-specific time and kriging standard deviation & Live arrival measurements, calibrated true-arrival distribution, or support beyond final perimeter.\\
Environmental alignment & Common grid values, categories, support metadata, and masks & Finer independent information than each native source or proven cutoff publication availability.\\
Probability predictor & Conditional next-hour transition estimates on its saved input and contact contract & Within-hour arrival density, an active flame front, spontaneous ignition, or physical propagation trajectories.\\
Residence assignment & Deterministic fuel-reference or explicit duration scenario & Measured flame geometry, a random duration law, crown/residual duration, or species-wide validated duration.\\
Occupancy bridge & Activity marginals under uniform arrival and constant duration & Learned arrival-time shape, existing-fire age, or a joint distribution of activity across cells.\\
Radiation operator & Incident mean flux, endpoint flux, dose, and conditional missing-source bounds & Heat variance, exceedance probability, absorbed flux, temperature, convection, terrain shielding, or exterior-source contribution.\\
\bottomrule
\end{longtable}

\subsection{Why expectation cannot determine a danger probability}
Let receiver flux be $F(x)=\sum_i k_i(x)Z_i$, with deterministic nonnegative
source-transfer coefficients $k_i(x)$. Marginal active probabilities determine
$\E[F(x)]=\sum_i k_i(x)p_i$. They do not determine
\begin{equation}
\Var(F(x))=\sum_i k_i(x)^2p_i(1-p_i)
+2\sum_{i<j}k_i(x)k_j(x)\Cov(Z_i,Z_j).
\end{equation}
Covariances and the full joint law are absent. For a concrete counterexample,
two sources of unit strength can each have marginal probability 0.4. If they
are perfectly correlated, total flux is 2 with probability 0.4 and 0 with
probability 0.6. If they are mutually exclusive, total flux is 1 with
probability 0.8 and 0 with probability 0.2. Both have mean 0.8, but
$\mathbb P(F>1.5)$ is 0.4 in the first case and zero in the second.
Thus a threshold applied to the expected field is a threshold on a mean
scenario output, not an estimated probability of harmful heat exposure.

\subsection{Conditions needed for an operational interpretation}
An operational evaluation would need a current fire-state product actually
available at each cutoff, feature availability and latency records, future
outcomes independent of reconstruction inputs, sufficient independent events,
and appropriate calibration of the intended endpoint. Active flames would
add event-specific arrival-age and duration evidence at a commensurate time
scale. Heat exposure would add fuel/consumption/radiation measurements and
validation for the intended receiver geometry and heat mechanism. None of
these missing quantities is inferred merely from a passing numerical test,
stored checkpoint, map image, or the algebraic timing bridge.

\appendix
\section{Source map and evidence conventions}
\label{app:sources}
Paths in this appendix are relative to the project root. The companion
\path{documents/wildfire_methods_source_audit.json} records SHA-256 hashes of
the inspected implementation and key saved records at document creation.
It also records current verification separately from saved historical reports.
The document is a description of that snapshot; later code edits require
checking these bindings before treating it as unchanged documentation.
\begin{longtable}{@{}P{0.27\linewidth}P{0.65\linewidth}@{}}
\toprule Method / evidence & Primary local location\\\midrule
\endfirsthead
\toprule Method / evidence & Primary local location\\\midrule
\endhead
Published reconstruction & \path{wildfire_burned_area_hourly/sources/zenodo_21788344/extracted/code/} and its retained \path{README.md}, \path{results/}, and original rasters.\\
Source provenance & \path{wildfire_burned_area_hourly/source_manifest.json} and \path{records/}.\\
Environmental acquisition & \path{wildfire_burned_area_hourly/environmental_data/acquire_environment.py}, \path{enrich_environment.py}, \path{audit_source_quality.py}, and \path{acquire_biome.py}.\\
Tensor and label contract & \path{wildfire_burned_area_hourly/pytorch_loader/dataset.py}.\\
Preparation and sampling & \path{wildfire_burned_area_hourly/probability_model/prepare_training.py} and \path{filter_contact.py}.\\
Exponential family & \path{wildfire_burned_area_hourly/probability_model/model.py}, \path{train.py}, and \path{inference.py}.\\
Signed direct family & \path{wildfire_burned_area_hourly/probability_model/direct_convolution/model.py}, \path{train.py}, \path{inference.py}, and \path{evaluate.py}.\\
Schema and scoring & \path{wildfire_burned_area_hourly/probability_model/input_schema.py}, \path{metrics.py}, and \path{broader_comparison.py}.\\
Saved kernel comparison & \path{wildfire_burned_area_hourly/probability_model/training_runs/20261007_1h_direct_convolution_signed/} and \path{20261007_1h_kernel_sweep/}.\\
Public model API & \path{wildfire_burned_area_hourly/api.py} and \path{api_example.py}.\\
Ecological loading & \path{expected_heat_flux/ecology.py}.\\
Residence scenarios & \path{expected_heat_flux/residence.py}, \path{sources.json}, and \path{documents/flame_residence_literature.md}.\\
Occupancy and flames & \path{expected_heat_flux/temporal.py} and \path{flames.py}.\\
Radiation and composition & \path{expected_heat_flux/model.py} and \path{pipeline.py}.\\
Saved heat demonstrations & \path{expected_heat_flux/outputs/integrated_demo/summary.json} and \path{active_flame_demo/summary.json}.\\
Current contract tests & \path{tests/loader/}, \path{tests/probability/}, \path{tests/heat/}, and \path{tests/test_api.py}.\\
\bottomrule
\end{longtable}

Claims labeled ``saved'' describe the preserved experiment or demonstration,
including its configured parameters and date. Statements about current
equations and contracts were checked against the implementation during
authorship. No probability models were refitted and no publisher study was
rerun for this document. Source-data integrity validation is distinct from
predictive or physical validation. In particular, the recorded full
environmental audit and its later biome extension explicitly distinguish
previously checked unchanged files from newly reread raster bands.

\section{Reproducibility and numerical checks}
The document itself is a standalone LaTeX source: all prose, mathematics,
tables, examples, and bibliography are contained in one \path{.tex} file.
It requires ordinary TeX packages, but no external figure, bibliography,
data, or code file to typeset. The companion PDF is an exported build of that
source. LaTeX compilation, resolution of cross-references/citations, final-log
diagnostics, and rendered-page layout are checked before delivery.

The one-cell calculation in Section~\ref{sec:worked} is checked through the
public temporal and radiation interfaces. The saved CPU inference example is
rerun to verify shape, schema, counts, and a known-probability range. Selected
residence and ecological-reader contract tests are run against the current
project source. Exact commands, counts, outputs, and source hashes are recorded
in the companion audit. These checks support the documented mathematical and
software behavior; they are not a new performance or field-calibration study.

\begin{thebibliography}{99}
\bibitem{toadeposit}
C. Choi, H. Seo, and J. Cha.
\emph{Wildfire time-of-arrival dataset and method code: variance-weighted
kriging of satellite active-fire detections, evaluated in the United States
and applied to South Korea}. Zenodo, record 21788344, 26 September 2026.
\url{https://doi.org/10.5281/zenodo.21788344}.
Companion article: \emph{Remote Sensing} 18(19), 3321 (2026),
\url{https://doi.org/10.3390/rs18193321}.
Record metadata and archived method code were inspected locally.

\bibitem{codetoa}
Choi, Seo, and Cha, retained dataset method code.
\path{preprocess.py}, \path{toa_kriging.py}, \path{produce_rasters.py},
\path{block_cv.py}, and \path{evaluate_sc.py} under the archived code
directory identified in Appendix~\ref{app:sources}. Units and source
conventions follow these archived files.

\bibitem{codeprob}
WildfireGuardian probability implementation and saved experiment records,
inspected 8 October 2026. Sources and records are located in
\path{wildfire_burned_area_hourly/probability_model/}; see
Appendix~\ref{app:sources} and the companion source-hash audit.

\bibitem{codeheat}
WildfireGuardian radiation, temporal, flame, and composed API implementation,
inspected 8 October 2026. \path{expected_heat_flux/model.py},
\path{temporal.py}, \path{pipeline.py}, and \path{flames.py}; saved
demonstrations are explicitly identified in the text.

\bibitem{heatrefs}
K.~J. Overholt.
\emph{Verification and Validation of Commonly Used Empirical Correlations
for Fire Scenarios}. NIST Special Publication 1169, March 2014,
Section~6.1, equation~(6.1), printed p.~39 (PDF page~53).
\url{https://doi.org/10.6028/NIST.SP.1169}.
\url{https://nvlpubs.nist.gov/nistpubs/SpecialPublications/NIST.SP.1169.pdf}.
The isotropic point-source law includes the receiver-orientation cosine.
The local radiation method integrates that differential source law over a
horizontal finite footprint; this is a package modeling extension.
A. Angelidis, \emph{Spherical Springs}, University of Otago Technical Report
OUCS-2004-01, 2004, Section~3.2, printed p.~2,
\url{https://cs.otago.ac.nz/research/publications/oucs-2004-01.pdf},
records the triangle solid-angle formula and attributes it to A. van Oosterom
and J. Strackee, \emph{The Solid Angle of a Plane Triangle}, IEEE Transactions
on Biomedical Engineering BME-30(2), 125--126 (1983),
\url{https://doi.org/10.1109/TBME.1983.325207}.
The accessible NIST and Otago PDFs were checked during preparation; the
original IEEE paper is an attribution from the accessible report, not a
claimed full-text inspection. The complete reference record is retained in
\path{expected_heat_flux/sources.json}.


\bibitem{coderesidence}
WildfireGuardian source and saved artifacts inspected on 8 October 2026.
\path{expected_heat_flux/residence.py},
\path{expected_heat_flux/ecology.py},
\path{expected_heat_flux/sources.json},
\path{tests/heat/test_residence.py},
\path{tests/heat/test_ecology.py},
\path{documents/flame_residence_literature.md}, and
\path{expected_heat_flux/outputs/active_flame_demo/summary.json}.
The separate $300\,\mathrm s$ integrated combustion scenario is recorded in
\path{expected_heat_flux/outputs/integrated_demo/summary.json}.

\bibitem{landfirefccs}
LANDFIRE.
\emph{Fuel Characteristic Classification System Fuelbeds}, official product
description and loading-table link.
\url{https://www.landfire.gov/fuel/fccs}.
Product description inspected online on 8 October 2026.
Reference loading table:
\url{https://www.landfire.gov/sites/default/files/CSV/LF_ConsumeLoadings.csv}.
The existing workflow uses archived inputs with recorded source epochs;
the current web page is not substituted for the archived raster release.

\bibitem{prichard2013}
S.~J. Prichard, D.~V. Sandberg, R.~D. Ottmar, E. Eberhardt, A. Andreu,
P. Eagle, and K. Swedin.
\emph{Fuel Characteristic Classification System version 3.0: technical
documentation}. USDA Forest Service, Pacific Northwest Research Station,
General Technical Report PNW-GTR-887, 2013.
\url{https://doi.org/10.2737/PNW-GTR-887}.
Official agency record and abstract:
\url{https://research.fs.usda.gov/treesearch/45283}.

\bibitem{anderson1969}
H.~E. Anderson.
\emph{Heat transfer and fire spread}. USDA Forest Service,
Research Paper INT-69, 1969, printed pp.~15--16.
\url{https://research.fs.usda.gov/download/treesearch/40489.pdf}.

\bibitem{andrews2018}
P.~L. Andrews.
\emph{The Rothermel surface fire spread model and associated developments:
A comprehensive explanation}. USDA Forest Service, Rocky Mountain Research
Station, General Technical Report RMRS-GTR-371, 2018.
\url{https://doi.org/10.2737/RMRS-GTR-371}.
Official PDF:
\url{https://research.fs.usda.gov/download/treesearch/55928.pdf}.
Surface-area weighting: Sections~3.2.3--3.2.5, printed pp.~21--24.
Residence relation and scope: Section~4.2, printed p.~27.
These sections were checked in the available official-PDF local snapshot;
the browsing service could not retrieve the large PDF in this audit.

\bibitem{scottburgan2005}
J.~H. Scott and R.~E. Burgan.
\emph{Standard fire behavior fuel models: a comprehensive set for use with
Rothermel's surface fire spread model}. USDA Forest Service, Rocky Mountain
Research Station, General Technical Report RMRS-GTR-153, 2005.
\url{https://doi.org/10.2737/RMRS-GTR-153}.
\url{https://research.fs.usda.gov/download/treesearch/9521.pdf}.
Component parameters: Table~7, printed p.~18. Common constants: p.~17.
Dynamic curing: p.~6. The source GS4 summary SAV differs from its table-derived
value; the implementation records and resolves this as explained in the text.
The relevant local official-PDF snapshot was inspected.

\bibitem{burrows1994}
N.~D. Burrows.
\emph{Experimental development of a fire management model for jarrah
(Eucalyptus marginata Donn ex Sm.) forest}.
PhD thesis, Australian National University, 1994.
Government-hosted original thesis:
\url{https://library.dbca.wa.gov.au/FullTextFiles/016094.pdf}.
Individual-particle methods: printed p.~67 (PDF page~76).
Implemented diameter equation: printed p.~69, equation~6-1 (PDF page~78).
These scanned pages were visually checked in the available local PDF;
the browsing service could not retrieve the large PDF in this audit.


\end{thebibliography}
\end{document}
